\chapter{Research Methodology}
\noindent{This study employed a multi-phase methodology to develop, calibrate, and validate a parallel hybrid agent-based metaheuristic optimization model for multi-route jeepney networks in Iligan City. The overall methodological framework integrates primary data collection through a structured behavioral survey, empirical derivation of travel cost parameters, construction of a multi-layered travel graph, simulation of passenger behavior through agent-based modeling, and optimization of route configurations using a metaheuristic algorithm. The present chapter describes each of these phases in detail, beginning with the design and execution of the commuter survey, which served as the principal source of behavioral data for the entire modeling pipeline.}

% 3.1. Survey Design and Data Collection

\section{Survey Design and Data Collection}
\label{sec:3.1 survey design and data collection}
\noindent
To ground the model in local commuter behavior rather than assumed or imported values, a structured behavioral survey was designed and administered to jeepney riders in Iligan City. The survey was hosted on Google Forms, conducted bilingually in English and Cebuano, and distributed through university networks, community platforms, and snowball referrals during May~\(2026\). A total of \(214\) valid responses were collected and retained for analysis. All respondents provided consent, and no duplicate usernames were identified during data cleaning.

The primary output of this phase was a set of empirically calibrated travel cost parameters expressed in \ac{EIVM}. These parameters use the canonical model terms \texttt{walk\_wt}, \texttt{ride\_wt}, \texttt{wait\_wt}, and \texttt{transfer\_wt}, together with the structural transition terms \texttt{direct\_wt} and \texttt{alight\_wt}. The final calibrated values were \texttt{walk\_wt} = 0.05630 \ac{EIVM}/m, \texttt{ride\_wt} = 0.00632 \ac{EIVM}/m, \texttt{wait\_wt} = 14.44 \ac{EIVM}/waiting event, \texttt{transfer\_wt} = 15.78 \ac{EIVM}/transfer event, \texttt{direct\_wt} = 0.00 \ac{EIVM}/transition, and \texttt{alight\_wt} = 0.00 \ac{EIVM}/transition. Distance-based parameters, namely \texttt{walk\_wt} and \texttt{ride\_wt}, are expressed in \ac{EIVM} per meter because they are multiplied by walking or riding distance in the \texttt{TravelGraph}. Event-based parameters, namely \texttt{wait\_wt} and \texttt{transfer\_wt}, are expressed in \ac{EIVM} per event because they are applied once when a passenger waits, boards, or transfers. The zero-cost structural parameters \texttt{direct\_wt} and \texttt{alight\_wt} are expressed in \ac{EIVM} per transition and are assigned zero to prevent double-counting costs already represented by walking, riding, waiting, or transferring.

These calibrated parameters are later implemented explicitly in the \texttt{TravelGraph} cost function described in Section \ref{sec:3.4 three layer graph}, where the graph distinguishes walking edges, jeepney ride edges, boarding or waiting transitions, transfer transitions, direct walking transitions, and alighting transitions. Within that graph, \texttt{walk\_wt} and \texttt{ride\_wt} convert physical edge distance into generalized travel cost, while \texttt{wait\_wt} and \texttt{transfer\_wt} add fixed behavioral penalties for waiting and transferring. This allows the \(A^{*}\) pathfinding algorithm to evaluate alternative passenger journeys using a unified cost scale that reflects not only geometric distance, but also locally measured commuter tolerance for walking, waiting, and route transfers. The construction of the \texttt{TravelGraph} and the behavior of the pathfinding algorithm are discussed in Section \ref{sec:3.4 three layer graph}; the present section explains how the survey was designed, collected, cleaned, and used to derive the parameters applied in that graph.

\subsection{Survey Rationale and Objectives}
\label{subsec:3.1.1 survey rationale and objectives}
A fundamental requirement of behaviorally grounded public transport modeling is that route cost functions reflect the actual preferences and tolerances of the commuter population being served, rather than generic or externally imported values \citep{ortuzar2011modelling}. In practice, commuters do not minimize travel time alone; they trade off different components of the journey, such as walking distance, waiting duration, in-vehicle time, and the inconvenience of transfers, against one another in ways that are specific to the local environment, infrastructure quality, and cultural context. Failure to capture these local trade-offs produces route assignments that are geometrically plausible but behaviorally invalid.

Iligan City presents a particularly constrained modeling context. Unlike highly documented metropolitan transit systems, no publicly available, peer-reviewed dataset existed at the time of this study that quantified commuter behavioral parameters specific to Iligan jeepney riders. National-level Philippine transport data do not disaggregate sufficiently to resolve city-level behavioral differences, and published values from other Philippine cities such as Metro Manila or Baguio City, cannot be assumed to transfer to Iligan without empirical verification, given differences in urban density, road network topology, and commuter demographics. This gap in local data made primary data collection indispensable.

The survey was accordingly designed to fulfill four specific objectives: (1) to characterize the typical trip profile of Iligan City jeepney commuters, including door-to-door duration, access and egress walking times, and habitual waiting durations; (2) to measure commuter walking tolerance, defined as the maximum walking distance a rider will accept before substituting an alternative mode; (3) to measure commuter waiting tolerance, defined as the duration of waiting a rider will accept in exchange for a faster jeepney service; and (4) to quantify transfer aversion, defined as the minimum time savings a commuter requires before accepting a route requiring a jeepney change in preference to a slower direct route. These four behavioral dimensions correspond directly to the four calibratable parameters of the travel cost model and constitute the complete behavioral input required for empirical graph calibration.

\subsection{Survey Structure}
\label{subsec:3.1.2 survey structure}
\noindent{The survey instrument was structured into six thematic sections: (A) Usual Trip, (B) Walking Tolerance, (C) Waiting Tolerance, (D) Transfer Behavior, (E) Demographics, and (F) Additional Comments; Administered bilingually in English and Cebuano. The bilingual format was adopted to ensure comprehension among the broadest possible cross-section of Iligan City commuters and to minimize response error attributable to language unfamiliarity. The survey was hosted on Google Forms and distributed electronically during May 2026. Participation was voluntary, and all respondents were required to provide informed consent prior to completing any substantive question, in compliance with the Data Privacy Act of 2012 (Republic Act No. 10173). Contact information for the research team and supervising faculty member was disclosed in the consent section to enable participants to raise data privacy concerns.}

\indent\textbf{Usual Trip (Questions A1–A8)} collected information on the respondent's habitual jeepney travel behavior. A1 asked about riding frequency, providing options ranging from daily use to never. A2 captured the primary trip purpose using a multi-select format. A3 and A4 recorded the approximate origin and destination of the respondent's most common trip using open-ended free-text fields. A5 measured total door-to-door trip duration using a categorical scale. A6 and A7 measured usual walking time to and from the jeepney stop, respectively, using five-category ordinal scales. A8 measured usual waiting time at the stop. The variables A6, A7, and A8 were used in Section \ref{subsec:3.1.4 parameter extraction and derivation} to establish baseline walking and waiting behavior.

\indent{\textbf{Walking Tolerance (Questions B1–B4)} measured the maximum walking distance a commuter will accept under different conditions. B1 and B2 presented respondents with a visual walking-time-ring map centered on Gaisano Main, Iligan City, and asked for maximum \ac{WTW} to a stop (access) and from a stop (egress), respectively. B3 posed a hypothetical scenario contrasting a nearer stop with a 15-minute headway against a more distant stop with a 5-minute headway, thereby eliciting a revealed-preference style walk-versus-wait trade-off. B4 was a conditional follow-up administered only to respondents who chose the farther stop in B3; it asked at what additional walking distance the respondent would switch back to the nearer stop. Together, B3 and B4 provided the empirical basis for deriving the walking cost multiplier.}

\indent{\textbf{Waiting Tolerance (Questions C1–C3)} assessed commuter sensitivity to waiting time. C1 identified the threshold at which a respondent would abandon waiting for a jeepney and instead board a tricycle or habal-habal. C2 presented a stated-preference scenario contrasting a jeepney departing immediately with a 30-minute in-vehicle travel time against a faster jeepney requiring a short additional wait, reducing in-vehicle time to 20 minutes; this elicited the wait-versus-ride trade-off used to derive the waiting cost multiplier. C3 was an attitudinal item assessing the effect of ambient heat on boarding decisions.}

\indent{\textbf{Transfer Behavior (Questions D1–D4)}  measured transfer frequency and transfer aversion. D1 established baseline transfer experience. D2 presented a stated-preference scenario comparing a direct 40-minute route against a transfer route with a 3-minute intermediate wait but only 30 minutes of total journey time, quantifying the choice between convenience and speed. D3 then asked, independently of the D2 scenario, how many minutes of time savings would be needed for the respondent to choose a transfer route over a direct alternative. The D3 response distribution, combined with a binary logistic regression model, provided the empirical basis for the transfer penalty derivation. D4 was a Likert-type matrix asking respondents to rate the effect of four contextual factors, namely covered waiting area, familiarity with the transfer location, rain, and being in a hurry, on their willingness to transfer.

\indent{\textbf{Demographics (Questions E1–E6)} recorded age group, gender, current occupational status, personal vehicle ownership, average daily jeepney expenditure, and the presence of any physical condition limiting walking capacity. Demographic data were used to characterize the respondent sample and to assess the representativeness of responses.}

\indent{\textbf{Additional Comments (Questions F1–F3)} was an optional section inviting open-ended feedback on jeepney routes, stops, and services, as well as a contact field for participation in a raffle incentive to encourage response completion.}

\subsection{Sampling Strategy and Sample Size}
\label{subsec:3.1.3 sampling strategy and sample size}
\indent\textbf{Target Population and Sampling Frame.} The target population for this study was defined as Iligan City commuters who use jeepneys for intra-city travel. Because no comprehensive registry of jeepney riders exists, and because jeepney ridership is distributed across an informal and unscheduled public transport network, a probability-based sampling frame could not be constructed. The study therefore adopted a purposive convenience sampling approach, with the survey disseminated through university networks, community platforms, and snowball referrals within Iligan City. This approach is appropriate for stated-preference transport studies in which the objective is not to enumerate the entire commuting population, but to obtain sufficient behavioral observations for estimating route-choice parameters, walking tolerance, waiting sensitivity, and transfer aversion (Louviere, Hensher, Swait, 2000).

\indent\textbf{Sample Size Justification.} The minimum sample size for population-level proportion estimation was first assessed using Cochran’s (1977) formula for an unknown population size:
\[
n_0 = \frac{Z^2 \cdot p \cdot q}{e^2}
\]

\indent{where \(Z = 1.96\) for a \(95\%\) confidence level, \(p = q = 0.50\) for maximum variance (a conservative assumption when the true population proportion is unknown), and 
\(e = 0.05\) for a margin of error of \(\pm 5\%\). This yields:}
\[
n_0 = \frac{(1.96)^2 \cdot (0.50) \cdot (0.50)}{(0.05)^2}
= \frac{3.8416 \cdot 0.25}{0.0025}
\approx 385
\]

\indent
The achieved sample size of \(N = 214\) valid responses is below the strict Cochran value required for population-level proportion estimation at \(\pm 5\%\) precision. However, this study does not use the survey to estimate citywide ridership proportions. Instead, the survey is used to calibrate behavioral parameters for a simulation model. Under this purpose, the adequacy of the sample is evaluated based on whether it provides sufficient observations for descriptive statistics, stated-preference trade-off analysis, and single-predictor logistic calibration.

For reference, if the collected \(N = 214\) were treated as a simple random sample under maximum variability, it would correspond approximately to a \(\pm 6.7\%\) margin of error:
\[
e = Z \sqrt{\frac{p(1-p)}{n}}
\]

\[
e = 1.96 \sqrt{\frac{0.50(0.50)}{214}}
\approx 0.0669
\]

This value is reported only as a descriptive benchmark, because the sampling design was non-probability-based. More importantly, the behavioral calibration tasks in this study are supported by the available response counts. For the transfer acceptance model in Section \ref{subsec:3.1.4 parameter extraction and derivation}, all 214 respondents contributed to the threshold-based transfer analysis. The threshold-to-choice expansion produced 856 binary stated-choice observations across four evaluated time-savings levels, with 346 acceptance observations and one predictor variable, time savings. This comfortably satisfies common logistic modeling guidance requiring at least 10 to 20 outcome events per predictor variable \citep{peduzzi1996simulation}. The sample size is therefore adequate for parameter extraction and simulation calibration, even though it is not claimed to provide strict population-level representation.

\indent\textbf{Data Yield and Cleaning.} A total of 214 survey responses were received. All respondents provided consent, and no duplicate usernames were identified after review. Therefore, no response was removed during cleaning, and the final analytic dataset comprised:
\[
N = 214 \text{ valid responses}
\]

All 214 responses were retained for descriptive analysis, while conditional variables were analyzed using only the relevant valid responses. For example, the B4 walk-wait switch-back question was evaluated only for respondents who selected Stop B in B3 and provided a valid switch-back threshold. Similarly, D3 finite transfer-saving thresholds were summarized separately from “Never” responses, while the full D3 distribution was retained in the logistic transfer acceptance model.

\indent\textbf{Respondent Profile.} The sample was strongly concentrated among young adult commuters. Respondents aged \(18\text{--}24\) accounted for \(199\) out of \(214\) responses, or \(93.0\%\), while the remaining \(15\) respondents, or \(7.0\%\), were aged \(25\text{--}34\). This age distribution is consistent with the university-centered dissemination strategy and the high dependence of students on jeepneys for daily mobility.

Occupational status further reflects this sampling channel. A total of \(200\) respondents, or \(93.5\%\), identified as students; \(10\) respondents, or \(4.7\%\), were employed in the private sector; and \(4\) respondents, or \(1.9\%\), were employed in the government sector. The gender distribution was \(112\) female respondents, or \(52.3\%\); \(98\) male respondents, or \(45.8\%\); and \(4\) respondents, or \(1.9\%\), who preferred not to disclose their gender.

In terms of ridership intensity, \(89\) respondents, or \(41.6\%\), reported riding jeepneys daily, while \(84\) respondents, or \(39.3\%\), reported riding several times per week. Together, \(173\) respondents, or \(80.8\%\) of the sample, were habitual or near-habitual jeepney users. This supports the construct validity of the sample because the majority of respondents had direct and repeated experience with jeepney waiting, walking access, boarding decisions, and transfer choices.

Because trip purpose was collected through a multi-select question, responses were not mutually exclusive. The most frequently reported trip purposes were shopping or errands with \(168\) respondents, or \(78.5\%\); school or university with \(165\) respondents, or \(77.1\%\); and going home with \(141\) respondents, or \(65.9\%\). These results indicate that although the sample is student-heavy, the reported jeepney trips covered a practical range of everyday urban travel purposes. Overall, the sample is appropriate for calibrating commuter behavior within the model, particularly for the walking, waiting, and transfer decisions that directly affect route choice in the travel graph.

\subsection{Parameter Extraction and Derivation}
\label{subsec:3.1.4 parameter extraction and derivation}
The travel cost model underlying the \texttt{TravelGraph} assigns generalized cost to each edge according to edge type and physical attributes. The total path cost is expressed in units of \textbf{\ac{EIVM}}. One \ac{EIVM} represents the perceived burden of one minute spent riding inside the jeepney. This unit is used because the model must compare walking distance, jeepney riding distance, waiting events, and transfer events within a single route-choice equation. Expressing all components in \ac{EIVM} allows the A$^{*}$ algorithm to compare alternative journeys using one consistent generalized cost scale \cite{ortuzar2011modelling}.

The total cost of a path is computed as:

\[
C(\text{path}) = \sum_{(u,v) \in \text{path}} \text{cost}(u,v)
\]

In implementation, the \texttt{TravelGraph} applies the calibrated parameters according to the edge type. Walking edges in the start-walk and end-walk layers are assigned cost as \texttt{walk\_wt} multiplied by edge length in meters. Ride edges in the jeepney layer are assigned cost as \texttt{ride\_wt} multiplied by ride-edge length in meters. Wait edges are assigned the fixed value \texttt{wait\_wt}, transfer edges are assigned the fixed value \texttt{transfer\_wt}, direct walking transitions are assigned \texttt{direct\_wt}, and alighting transitions are assigned \texttt{alight\_wt}. Therefore, \texttt{walk\_wt} and \texttt{ride\_wt} are distance-based coefficients measured in \ac{EIVM}/m, while \texttt{wait\_wt} and \texttt{transfer\_wt} are event-based penalties measured in \ac{EIVM}/event. The structural transition parameters \texttt{direct\_wt} and \texttt{alight\_wt} are measured in \ac{EIVM}/transition and are assigned zero cost to avoid double-counting costs already captured by walking, riding, waiting, or transferring.

\noindent\textbf{Data Preprocessing and Statistical Treatment.}
The survey used categorical and ordinal response choices. To make the responses usable for computation, time categories were converted into numeric values using midpoint or threshold coding. Midpoint coding was used for ordinary duration ranges because respondents selected an interval rather than an exact value. For example, ``2--5 minutes'' was coded as 3.5 minutes, and ``5--10 minutes'' was coded as 7.5 minutes. Open-ended upper categories were conservatively capped using a lower-bound extension: ``10 minutes or more'' was coded as 12 minutes for walking tolerance, while ``15 minutes or more'' was coded as 17.5 minutes for waiting tolerance. This prevents extreme assumptions from inflating the derived weights while still acknowledging that the response exceeds the previous interval. This treatment is consistent with stated-preference survey analysis, where categorical responses must often be converted into analyzable continuous values while avoiding unsupported assumptions about unobserved extremes \cite{hensher2015applied}.

For threshold-type questions, endpoint coding was used instead of midpoint coding because the response represents a maximum accepted value or a minimum required condition rather than an average duration inside an interval. For C2, ``Up to 5 minutes of waiting'' was coded as 5 minutes because the respondent indicates a maximum accepted wait. For D3, ``20 minutes or more of savings'' was coded conservatively as 20 minutes because it represents the lower bound of the required saving. For B4, ``Beyond 4,'' ``Beyond 6,'' and ``Beyond 8 minutes of walking'' were coded as 4, 6, and 8 minutes, respectively, as lower-bound switch-back thresholds.

The following table summarizes the recalculated descriptive statistics for the variables used in parameter derivation.

\input{chap3/tables/table_3.1}

The 85th percentile of pooled \ac{WTW} was selected as the behavioral access-radius reference because the objective at this stage was to define a practical system boundary for most users, not to estimate an average walking cost. A mean or median walking distance would describe central tendency, but it would exclude a substantial share of commuters who are still realistically willing to walk farther. \cite{elgeneidy2014walking} supports the use of high-percentile walking catchments for transit access analysis, noting that ``the 75th and 85th percentile buffers, which more accurately represent walking area for most users'' provide a more realistic service-area boundary than a single fixed average. Therefore, the 85th percentile was used here to represent a high-coverage but still empirically observed walking-access threshold.

Using a baseline walking speed of 1.2 m/s, equivalent to 72 m/min, the 12-minute 85th percentile walking tolerance corresponds to:

\[
12 \,\text{min} \times 60 \,\text{s/min} \times 1.2 \,\text{m/s} = 864 \,\text{meters}
\]

Thus, 864 meters was retained as the behavioral access-radius reference for the model. This means that approximately 85\% of the sampled walking-tolerance responses fall within a 12-minute, or 864-meter, walking threshold. The baseline walking speed of 1.2 m/s is consistent with pedestrian speed literature: \cite{katsaros2024} provide a systematic review of pedestrian walking speed analysis, while \cite{hassouna2020evaluation} discusses pedestrian walking speed patterns at crosswalks and supports the continued relevance of 1.2 m/s as a typical pedestrian design-speed reference. This access-radius value was used only as a behavioral reference for stop accessibility interpretation. It was not used directly as \texttt{walk\_wt} because \texttt{walk\_wt} must represent marginal generalized cost per meter, not a maximum acceptable walking distance.

\noindent\textbf{Rationale for Using Mean, Median, Percentile, and Logistic Thresholds.}
Different statistics were used because each parameter represents a different behavioral concept. A single statistic cannot validly represent all components.

The mean was used for A8 usual waiting time because \texttt{wait\_wt} represents the expected cost of a typical waiting event. Since simulation cost accumulates across many passengers and boarding episodes, the average waiting burden is the appropriate event-level baseline.

The median was used for the C2 wait-versus-ride trade-off because C2 is a stated threshold question with a skewed response distribution and an open-ended upper category. The median identifies the typical break-even respondent and is less distorted by respondents who would not wait at all or respondents willing to wait 15 minutes or more. In this dataset, the median accepted wait for a 10-minute faster ride was 5 minutes, producing a clear waiting multiplier of 2.0 \ac{EIVM}/min waiting.

The mean was used for the valid B4 switch-back responses because \texttt{walk\_wt} is a marginal distance-based coefficient applied continuously across walking edges. For a per-meter edge cost, the expected additional walking tolerance among respondents who accepted the walk-for-wait trade-off is more appropriate than the median alone. Using only the median would ignore the distribution of respondents willing to continue walking farther and would over-penalize walking relative to the observed trade-off behavior.

The 85th percentile was used only for the walking-access reference boundary because that value describes how far the model may reasonably treat stops as accessible for most users. It was not used to compute \texttt{walk\_wt}, because a percentile boundary describes coverage, while a cost weight describes marginal burden.

A logistic threshold was used for \texttt{transfer\_wt} because transfer acceptance is a binary choice process. At a given time saving, a commuter either accepts or rejects the transfer. The D3 response also includes ``Never'' responses, which cannot be assigned a finite numerical threshold without imposing an arbitrary value. Binary logistic modeling allows these responses to be retained as non-acceptance observations within the observed savings range, producing a defensible behavioral break-even point \citep{benakiva1985discrete,train2009discrete}.

\noindent\textbf{Derivation of \texttt{ride\_wt}.}
\texttt{ride\_wt} converts jeepney travel distance in meters into \ac{EIVM}. Since the \ac{EIVM} scale uses in-vehicle time as its reference unit, the riding coefficient is based on the expected number of in-vehicle minutes required to travel one meter by jeepney.

A baseline jeepney operating speed of 9.5 km/h was adopted from Philippine jeepney operating-speed literature \cite{ranosa2021}. Converting this speed into meters per minute gives:

\[
v_{\text{jeep}} = 9.5 \text{ km/h} \times \frac{1000 \text{ m}}{1 \text{ km}} \times \frac{1 \text{ h}}{60 \text{ min}}
\]
\[
v_{\text{jeep}} = 158.33 \text{ m/min}
\]

Since one minute of jeepney riding is equal to one \ac{EIVM}, the riding weight is:

\[
\texttt{ride\_wt} = \frac{1 \text{ \ac{EIVM}/min}}{158.33 \text{ m/min}}
\]
\[
\texttt{ride\_wt} \approx 0.00632 \text{ \ac{EIVM}/m}
\]

This means that each meter of jeepney travel adds approximately 0.00632 \ac{EIVM} to the generalized path cost before any congestion multiplier is applied.

\noindent\textbf{Derivation of \texttt{wait\_wt}.}
\texttt{wait\_wt} is an event-based penalty applied when a passenger boards or waits for a jeepney. Its unit is \ac{EIVM}/waiting event. It is not multiplied by distance because waiting occurs as a discrete episode rather than as movement along a road segment.

The derivation used two survey variables. First, A8 measured respondents' usual waiting time. After midpoint coding, the mean usual waiting time was:

\[
\bar{W}_{\text{survey}} = 7.22 \text{ min/waiting event}
\]

Second, C2 measured how much waiting respondents would accept to save 10 minutes of jeepney in-vehicle time. The median accepted wait was 5 minutes. Therefore, the waiting-time multiplier is:

\[
\theta_{\text{wait}} = \frac{10 \text{ \ac{EIVM}}}{5 \text{ min waiting}}
\]
\[
\theta_{\text{wait}} = 2.0 \text{ \ac{EIVM}/min waiting}
\]

This means that one minute of waiting is perceived as equivalent to two minutes of in-vehicle jeepney riding. This is consistent with public transport value-of-time literature, which generally finds waiting time to be valued more heavily than in-vehicle time \citep{wardman2004public,itfoecd2014valuing}.

The waiting event penalty is then:

\[
\texttt{wait\_wt} = \bar{W}_{\text{survey}} \times \theta_{\text{wait}}
\]
\[
\texttt{wait\_wt} = 7.22 \text{ min/event} \times 2.0 \text{ \ac{EIVM}/min}
\]
\[
\texttt{wait\_wt} = 14.44 \text{ \ac{EIVM}/waiting event}
\]

Thus, each boarding or waiting episode adds 14.44 \ac{EIVM} to the path cost. This is equivalent to the perceived burden of approximately 14.44 minutes of in-vehicle jeepney travel.

\noindent\textbf{Derivation of \texttt{walk\_wt}.}
\texttt{walk\_wt} converts walking distance into \ac{EIVM}. Its unit is \ac{EIVM}/m because walking edges are costed according to distance:

\[
C_{\text{walk}} = d \cdot \beta_{\text{walk}}
\]

where $d$ is walking distance in meters and $\beta_{\text{walk}}$ is \texttt{walk\_wt}.

A baseline walking speed of 1.2 m/s was used, equivalent to:

\[
v_{\text{walk}} = 1.2 \text{ m/s} \times 60 \text{ s/min}
\]
\[
v_{\text{walk}} = 72 \text{ m/min}
\]

The walking multiplier was derived from the B3/B4 walk-wait trade-off. In B3, respondents chose between Stop A, requiring a 2-minute walk and 15-minute wait, and Stop B, requiring a 6-minute walk and 5-minute wait. Choosing Stop B therefore required 4 additional minutes of walking but saved 10 minutes of waiting. Out of 214 respondents, 143 selected Stop B, representing 66.8\% of the sample. This indicates that most respondents were willing to walk farther if doing so substantially reduced waiting time.

The \ac{EIVM} benefit of saving 10 minutes of waiting is:

\[
10 \text{ min waiting} \times 2.0 \text{ \ac{EIVM}/min waiting} = 20 \text{ \ac{EIVM}}
\]

B4 then asked Stop B respondents when they would switch back to Stop A. Among respondents who chose Stop B and gave a valid B4 switch-back threshold, $N = 137$. The mean switch-back walking threshold to Stop B was 6.93 minutes. Since Stop A required a 2-minute walk, the mean additional walking tolerance relative to Stop A was:

\[
6.93 \text{ min} - 2.00 \text{ min} = 4.93 \text{ additional walking minutes}
\]

The walking multiplier is therefore:

\[
\theta_{\text{walk}} = \frac{20 \text{ \ac{EIVM}}}{4.93 \text{ min walking}}
\]
\[
\theta_{\text{walk}} \approx 4.05 \text{ \ac{EIVM}/min walking}
\]

Converting this from \ac{EIVM} per walking minute into \ac{EIVM} per meter:

\[
\texttt{walk\_wt} = \frac{\theta_{\text{walk}}}{v_{\text{walk}}}
\]
\[
\texttt{walk\_wt} = \frac{4.05 \text{ \ac{EIVM}/min}}{72 \text{ m/min}}
\]
\[
\texttt{walk\_wt} \approx 0.05630 \text{ \ac{EIVM}/m}
\]

Thus, every meter walked adds approximately 0.05630 \ac{EIVM} to the generalized path cost. This value is greater than \texttt{ride\_wt} because walking is slower than riding and also carries physical and environmental burdens such as heat exposure, rain exposure, sidewalk conditions, and fatigue. The result is consistent with transport modeling literature, which treats walking time as more burdensome than in-vehicle time \citep{wardman2004public,ortuzar2011modelling}.

\noindent\textbf{Derivation of \texttt{transfer\_wt}.}
\texttt{transfer\_wt} is a fixed event penalty applied each time a passenger transfers from one jeepney route to another. Its unit is \ac{EIVM}/transfer event because it is applied once per transfer occurrence, rather than multiplied by distance. In the \texttt{TravelGraph}, this corresponds to the cost of a transfer transition from the end-walk layer back into the ride layer. This formulation separates transfer inconvenience from the measurable walking and waiting components of the trip.

This separation is academically necessary because transfer disutility is not fully explained by additional walking time or waiting time alone. \cite{garciamartinez2018transfer} explicitly distinguish the pure transfer penalty from other transfer-related components such as access time, waiting time, and in-vehicle time. Their discrete-choice analysis found that the pure transfer penalty can be equivalent to approximately 15.2 to 17.7 additional in-vehicle minutes, meaning that passengers may still prefer a longer direct trip over a faster transfer trip even when the transfer does not add extra walking or waiting. This supports the decision to model \texttt{transfer\_wt} as an independent fixed behavioral penalty rather than absorbing it into \texttt{walk\_wt} or \texttt{wait\_wt}.

The survey provided direct local evidence of transfer aversion. In D2, only 63 out of 214 respondents, or 29.4\%, selected the transfer route even though it saved 10 minutes compared with the direct route. In D3, 77 respondents, or 36.0\%, stated that they would never choose a transfer route over a direct route regardless of time savings. Among respondents who gave finite thresholds, 35 required 5 minutes of savings, 44 required 10 minutes, 16 required 15 minutes, and 42 required 20 minutes or more. These results indicate that transfer acceptance behaves as a threshold-based binary choice: at a given amount of time savings, a commuter either accepts or rejects the transfer route.

To convert this distribution into a transfer penalty, D3 responses were transformed into a binary stated-choice structure. Four candidate time-savings levels were evaluated:

\[
S \in \{5,10,15,20\}
\]

where $S$ is measured in minutes of in-vehicle time saved by choosing the transfer route instead of the direct route. For each respondent and each value of $S$, transfer acceptance was coded as:

\[
y_i =
\begin{cases}
1, & \text{if the respondent's required savings threshold is } \le S, \\
0, & \text{otherwise.}
\end{cases}
\]

Respondents who selected ``Never'' were coded as 0 across all four observed savings levels. This retained their stated preference without assigning an arbitrary numerical threshold beyond the survey range. The transformation produced:

\[
214 \text{ respondents} \times 4 \text{ savings levels} = 856 \text{ binary observations}
\]

A binary logistic model was then fitted using maximum likelihood estimation, following the discrete-choice modeling framework commonly used for transport stated-preference analysis \cite{benakiva1985discrete,train2009discrete}. The model is defined as:

\[
P(\text{accept transfer}) = \frac{1}{1+e^{-(\alpha+\beta S)}}
\]

where $P(\text{accept transfer})$ is the predicted probability of accepting the transfer route, $S$ is the time saved by the transfer route in minutes, $\alpha$ is the dimensionless intercept, and $\beta$ is the marginal increase in log-odds per additional minute of time savings, with unit $\text{min}^{-1}$.

The expression $\alpha+\beta S$ represents the logit model's linear utility index, or the log-odds of accepting the transfer route at a given level of time savings. In this framework, the parameters $\alpha$ and $\beta$ must be estimated from the survey data. Assuming that each respondent's choice is independent, the joint probability of observing the exact sequence of 856 ``accept'' and ``reject'' decisions in the dataset is the product of the individual probabilities for each observation. This forms the likelihood function:

\[
L(\alpha,\beta)=\prod_{i=1}^{856} P_i^{y_i}(1-P_i)^{1-y_i}
\]

Here, if an observation is an acceptance ($y_i=1$), the term simplifies to $P_i$. If it is a rejection ($y_i=0$), it simplifies to $(1-P_i)$.

However, multiplying 856 fractions together results in an infinitesimally small number, making optimization computationally difficult due to numerical underflow. Because the natural logarithm is a strictly increasing function, the values of $\alpha$ and $\beta$ that maximize the likelihood function will also maximize its logarithm. Taking the natural logarithm converts the mathematically unstable product into a much simpler, stable summation known as the log-likelihood:

\[
\ell(\alpha,\beta)=\sum_{i=1}^{856} \left[y_i \ln(P_i)+(1-y_i)\ln(1-P_i)\right]
\]

Because there is no simple algebraic way to solve this equation, the maximum likelihood procedure uses an iterative computational algorithm. The algorithm tests successive values of $\alpha$ and $\beta$, evaluating the log-likelihood score at each step, until it finds the exact parameter combination that maximizes the probability of observing the actual survey responses. This produces the logistic S-curve that most accurately fits the observed transfer acceptance pattern (see Figure \ref{fig:transfer_logistic}). The optimized coefficients that achieved this maximum likelihood were:

\[
\alpha=-2.1242
\]
\[
\beta=0.1346 \text{ min}^{-1}
\]

Visually, these parameters define the shape of the probability curve. The intercept $\alpha$ establishes the baseline negative utility: when no time is saved ($S=0$), the probability of accepting a transfer is very low, reflecting inherent transfer aversion. The coefficient $\beta$ determines the steepness of the curve: each additional minute of savings increases the log-odds of accepting the transfer by 0.1346, pulling the probability upward.

\begin{figure}[htbp]
    \centering
    \includegraphics[width=0.8\textwidth]{chap3/figures/logistic_curve.png}
    \caption{Fitted logistic curve showing the probability of accepting a transfer route versus in-vehicle time savings.}
    \label{fig:transfer_logistic}
\end{figure}

The transfer penalty was defined as the 50\% predicted acceptance threshold. This threshold is not arbitrary. In binary logit theory, a 50\% probability occurs when two alternatives have equal systematic utility \cite{train2009discrete}. In this study, $P=0.50$ therefore represents the behavioral break-even point at which the transfer route and the direct route are equally attractive to the modeled commuter population.

At $P=0.50$, the log-odds are zero:

\[
\ln\left(\frac{P}{1-P}\right)=\ln\left(\frac{0.50}{0.50}\right)=\ln(1)=0
\]

Therefore, the linear utility index must equal zero at the break-even time saving ($S^*$):

\[
\alpha+\beta S^*=0
\]

Solving for $S^*$:

\[
S^*=\frac{-\alpha}{\beta}
\]
\[
S^*=\frac{-(-2.1242)}{0.1346}
\]
\[
S^* \approx 15.78 \text{ min}
\]

Because the time-saving threshold is expressed in minutes of in-vehicle travel saved, and one in-vehicle minute is equal to one \ac{EIVM}, the pure transfer penalty is isolated as:

\[
\texttt{transfer\_wt}=15.78 \text{ \ac{EIVM}/transfer event}
\]

This value means that the model treats one transfer as equivalent to approximately 15.78 minutes of in-vehicle jeepney travel. In practical terms, a transfer route must save about 15.78 minutes before the typical modeled commuter becomes equally likely to accept the transfer instead of remaining on a direct route. The calibrated value is therefore locally derived from Iligan commuter survey data, while also being externally comparable to established public transport literature on pure transfer penalties, particularly the 15.2--17.7 in-vehicle-minute equivalence reported by \citep{garciamartinez2018transfer} and the broader international transfer-penalty equivalency discussion in \citep{jara-diaz2022}.

\noindent\textbf{Justification of \texttt{direct\_wt} and \texttt{alight\_wt}.}
\texttt{direct\_wt} and \texttt{alight\_wt} are structural transition weights, not independent behavioral penalties. Both are assigned a value of 0.00 \text{ \ac{EIVM}/transition}

This does not mean that direct walking or alighting has no role in the journey. Rather, it means that the transition edge itself does not add a separate cost beyond the travel burden already charged by the appropriate walking, riding, waiting, or transfer component.

This treatment follows the principle of generalized cost decomposition, where each component of travel burden should be counted once through its corresponding term \citep{ortuzar2011modelling}. It is also consistent with transit assignment models that represent public transport journeys as networks of distinct access, ride, alight, and transfer links, where each link type carries cost only when it represents an actual independent impedance \citep{spiess1989optimal}. Assigning positive cost to every transition would risk double-counting the same travel burden.

\noindent\textbf{Summary of Calibrated Parameters.}
The final calibrated travel cost parameters are summarized below.
\input{chap3/tables/table_3.2}

\section{Spatial Data Preparation and Network Construction}
\label{sec:3.2 CityGraph}

\indent {The spatial foundation of the proposed system is built from an \ac{OSM}-derived road network of Iligan City, converted into a directed graph representation suitable for route generation, passenger pathfinding, and simulation. This stage prepares the physical road skeleton on which the succeeding demand model, route generator, and passenger simulation operate. Rather than treating road data only as a map visualization layer, the study converts the city network into a computable graph composed of nodes, directed road segments, road-type eligibility flags, and adjacency relations. This graph is referred to as the \texttt{CityGraph}, and it serves as the base infrastructure object for all later route construction processes.}

\indent {Within the architecture, the \texttt{CityGraph} is treated as the Layer 0 Infrastructure Layer. This layer contains the physical street network before it is transformed into passenger-state layers in the later \texttt{TravelGraph} construction. The Layer 0 designation is important because it separates the extracted road infrastructure from the later pedestrian, transit, and destination states. Thus, the \texttt{CityGraph} remains the underlying spatial reference, while the \texttt{TravelGraph} later uses it to generate passenger \texttt{Journey} objects. In this study, a \texttt{Journey} refers to the ordered sequence of directed edges produced for a passenger \ac{OD} request. It represents the complete path the passenger is expected to follow through the network, but the detailed multi-layer mechanics of that path are discussed later in Section \ref{sec:3.4 three layer graph}.}

\indent {The \texttt{TravelGraph} is introduced here only as the structure that eventually consumes the \texttt{CityGraph} and the generated route system to compute passenger \texttt{Journey} objects. Its role is to translate infrastructure and route data into pathfinding-ready passenger movement. The \texttt{CityGraph}, however, is the main concern of this section. It defines which physical roads exist, which of those roads are eligible for jeepney route construction, how shortest paths are computed between candidate route points, and how raw geographic coordinates are snapped to valid topological nodes. This order is necessary because the route generator cannot construct valid jeepney lines until the physical network has first been cleaned, simplified, and classified.}

\indent {The cost structure also preserves the distinction between physical movement and structural transitions. Later in the \texttt{TravelGraph}, \texttt{direct\_wt} and \texttt{alight\_wt} are assigned zero cost because they function as state-change connectors rather than independent sources of travel burden. A direct connector allows a passenger path to continue as a non-transit walking alternative without adding an artificial penalty, while an alighting connector changes the passenger state from riding to post-ride walking without duplicating costs already represented by in-vehicle travel and walking edges. This follows the logic of generalized transit network representations, where certain connector arcs exist to maintain valid passenger movement states rather than to introduce additional travel time penalties \cite{spiess1989optimal}. Therefore, assigning these transition weights to zero does not remove travel cost from the \texttt{Journey}; it prevents the same physical movement from being counted twice.}

\subsection{OSMnx Extraction and Graph Simplification}

\indent {The \texttt{CityGraph} is a persistent infrastructure object that converts \ac{OSM}-derived road data into a typed graph representation usable by pathfinding, route generation, visualization, and later \texttt{TravelGraph} construction. Its internal representation is built from three core object types: \texttt{Node}, \texttt{DirEdge}, and \texttt{CityGraph}.}

\indent {\textbf{The Node Object.} A \texttt{Node} is the atomic spatial unit of the system. It stores a unique identifier, longitude, latitude, and layer designation. Coordinates are consistently represented in \texttt{(lon, lat)} order, and coordinate validation ensures that longitude values fall within $[-180, 180]$ and latitude values fall within $[-90, 90]$. In the \texttt{CityGraph}, every infrastructure node is assigned to Layer 0. This means that the physical node representing a road intersection remains distinct from any later node at the same coordinate that may be created for passenger walking, riding, or destination states.}

\indent {All distance computations involving nodes use the Haversine formula rather than planar Euclidean distance. This is necessary because road-network coordinates are expressed in longitude and latitude, and Haversine distance better approximates surface distance on the Earth for geographic coordinate pairs. In this study, those distances are stored and interpreted in meters.}

\indent {\textbf{The DirEdge Object.} A \texttt{DirEdge} is a directed connection from one \texttt{Node} to another. It stores its start node, end node, drivable status, traversal weight, unique identifier, outgoing edge list, and precomputed geographic length. The outgoing edge list, stored as \texttt{next\_edges}, is central to graph traversal because it determines which directed edges can be reached after completing the current edge.}

\indent {In the \texttt{CityGraph}, the main concern is the Layer 0 to Layer 0 road edge. The same \texttt{DirEdge} abstraction is later reused by the \texttt{TravelGraph} for additional movement and transition types, but those inter-layer semantics are discussed in Section \ref{sec:3.4 three layer graph}. For Section \ref{sec:3.2 CityGraph}, the important point is that physical road segments are represented as directed objects so that shortest-path algorithms can traverse the road network through explicit start-to-end movement.}

\indent {Bidirectional road movement is represented by creating two opposing directed edges for each retained road segment, one in each direction. This avoids the need to maintain separate undirected-edge semantics inside a directed pathfinding system. The operational restriction on which of these edges may be used by \ac{PUJ} routes is not imposed at this bidirectional-rendering step; it is handled separately by the \ac{PUJ} arterial function described in Section~3.2.2.}

\indent {\textbf{The CityGraph Object.} The \texttt{CityGraph} class assembles \texttt{Node} and \texttt{DirEdge} objects into a coherent road-network object. It stores the graph name, infrastructure nodes, directed infrastructure edges, optional landmarks, outgoing-edge indexes, and cache information. It is initialized using a geographic bounding box, an optional local \ac{PBF} file path, a \texttt{use\_api} flag, and cache settings:}

\indent {The preferred extraction pathway uses a locally stored \ac{OSM} \ac{PBF} file, or \texttt{.osm.pbf} file. This \ac{PBF} extract is obtained independently from the BBBike Extract Service, which allows a custom geographic bounding box to be selected for the study area. This workflow is preferred over immediate \ac{osmnx} \ac{API} extraction because it reduces repeated network requests, allows offline processing, and ensures that repeated optimization runs are based on the same geographic extract. The bounding box follows the project convention \texttt{(min\_lat, max\_lat, min\_lon, max\_lon)}, while all internal graph coordinates follow \texttt{(lon, lat)}.}

\indent {The local \ac{PBF} file is parsed using \texttt{pyrosm}. The implementation clips the extract to the Iligan City bounding box, extracts the driving network, and converts the resulting nodes and edges into a NetworkX \texttt{MultiDiGraph}. If the \ac{PBF} file is unavailable, the implementation can fall back to \ac{osmnx} extraction through the \ac{OSM} Overpass \ac{API} by setting \texttt{use\_api=True}. \ac{osmnx} is appropriate for this fallback because it provides established methods for downloading, constructing, simplifying, and analyzing street networks from \ac{OSM} data \citep{boeing2017osmnx, boeing2025}. However, the local \ac{PBF} workflow remains the primary pathway for reproducibility and efficiency.}

\indent {After extraction, the raw road graph is serialized into a binary pickle cache. The cache filename is generated using an \ac{MD5} hash computed from the cache prefix and bounding-box tuple. In this context, \ac{MD5} is not used as a security mechanism; it is used only as a deterministic hash function for generating compact and repeatable cache keys. This allows the same geographic bounds and cache prefix to map to the same cached road graph across repeated runs. \ac{MD5} was originally introduced as a message-digest algorithm for producing fixed-length identifiers from arbitrary input strings \citep{rivest1992md5}. Here, that property is used only for reproducible file naming. On subsequent runs, the cached graph is loaded directly, avoiding repeated \ac{PBF} parsing or \ac{API} extraction during the many initialization cycles required by metaheuristic route optimization.}

\indent {Graph construction proceeds in two phases. In the node construction phase, the implementation reads coordinate attributes from the extracted graph. \ac{osmnx}-style graphs store longitude and latitude as \texttt{x} and \texttt{y}, while \texttt{pyrosm} may expose equivalent \texttt{lon} and \texttt{lat} attributes. The \texttt{CityGraph} handles both conventions and creates a Layer 0 \texttt{Node} for every valid coordinate pair. Nodes without valid coordinates are skipped.}

\indent {In the edge construction phase, each retained \ac{OSM} road segment is converted into one pair of opposing \texttt{DirEdge} objects. Several simplification rules are applied before an edge is admitted into the graph. Self-loops, defined as edges whose start and end node identifiers are identical, are discarded. Geometrically degenerate edges, defined as edges whose start and end nodes have different \ac{OSM} identifiers but identical geographic coordinates, are also discarded. Duplicate undirected pairs are filtered using a \texttt{seen\_pairs} set so that repeated \ac{OSM} geometries do not create redundant road links. For each retained segment, two directed edges are created, one in each direction. The \texttt{\_is\_arterial()} method then evaluates the segment's \ac{OSM} \texttt{highway} tag and assigns the edge-level \texttt{is\_drivable} flag. The road-type rules behind this flag are detailed in Section \ref{subsec:3.2.2 puj arterial func}.}

\indent {The object-level validation rules for \texttt{Node} and \texttt{DirEdge}
construction are provided in Appendix~\ref{appendix:node-diredge-implementation}.
The appendix contains the implementation-level safeguards for coordinate validation,
layer validation, coordinate immutability, same-layer movement restrictions, and
inter-layer transition restrictions. In the methodology, these implementation details
are summarized as the graph-construction logic shown in Algorithm~\ref{alg:validated-node-diredge-construction}.}

\input{chap3/algorithms/alg_node_diredge_construction}

\indent {After node and edge construction, the graph is stitched. Stitching populates each edge's \texttt{next\_edges} list with the outgoing edges reachable from its terminal node. Rather than repeatedly scanning the full edge list, the helper function uses coordinate-and-layer keys to connect compatible edges efficiently. This produces the adjacency structure used by \texttt{CityGraph.find\_shortest\_path()}, which executes \(A^{*}\) pathfinding over the drivable subset of the infrastructure graph.}

\subsection{PUJ Arterial Function}
\label{subsec:3.2.2 puj arterial func}
\indent {Not all roads in the \texttt{CityGraph} are eligible for jeepney route generation. Although the \texttt{CityGraph} preserves the broader street network for spatial completeness, the \texttt{RouteGenerator} must be restricted to roads that plausibly function as \ac{PUJ} corridors. Without this restriction, the metaheuristic would be allowed to search through residential dead-ends, service alleys, pedestrian paths, and other low-access road classes that are not operationally suitable for jeepney routes. In transit network design, constraining the candidate search space is necessary because unfiltered networks can create combinatorial explosion and waste computation on infeasible or low-quality route alternatives \cite{iliopoulou2019metaheuristics}.}

\indent {The arterial restriction is implemented through the \texttt{is\_drivable} Boolean flag assigned to each \texttt{DirEdge}. In this study, \texttt{is\_drivable=True} does not mean that every vehicle can physically pass through the road. Instead, it means that the road segment is eligible for \ac{PUJ} route construction. Edges marked \texttt{is\_drivable=True} form the arterial skeleton used by the \texttt{RouteGenerator}, while edges marked \texttt{is\_drivable=False} remain available in the broader infrastructure graph for non-route spatial representation and later passenger walking access.}

\indent {The operational basis for this road-type taxonomy is supported by the Philippine paratransit hierarchy described by ~\cite{guillen2013}. Their study identifies \ac{PUJ}s as major public transport modes for longer urban trips, while tricycles and pedicabs serve shorter access trips that connect passengers to \ac{PUJ} routes. In this study, that hierarchy is used together with the fixed-route nature of \ac{PUJ} operation to justify restricting route generation to arterial and connector corridors rather than allowing generated routes to penetrate residential, service, pedestrian, or access-grade roads.}

\indent {The full road-type inclusion and exclusion taxonomy used by the \ac{PUJ} arterial function is summarized in Table~\ref{tab:puj-road-type-taxonomy}.}

\input{chap3/tables/table_3.3}

\subsection{Route Construction and Validation Pipeline}
\label{subsec:3.2.3 route construction and validation pipeline}
After the CityGraph has been constructed and its \ac{PUJ}-eligible arterial skeleton has been identified, candidate jeepney routes are generated through the RouteGenerator. The generator accepts a CityGraph instance and a DirectDemandSampler as inputs. This design makes route generation demand-responsive rather than uniformly random, because waypoint selection is biased toward nodes with higher estimated passenger relevance, as described later in Section \ref{sec:3.3 ddm}.

The route construction pipeline proceeds through four main stages: demand-weighted waypoint sampling, shortest-path connection and loop closure, Layer~2 promotion, and structural validation.

\textbf{Waypoint Sampling.}
For each candidate route, the RouteGenerator samples $n$ waypoints from the CityGraph using the probability distribution produced by the DirectDemandSampler. Sampling is constrained by \texttt{only\_drivable=True}, meaning that every selected waypoint must lie on the arterial skeleton defined by the \ac{PUJ} arterial function. This prevents route construction from beginning or terminating inside non-\ac{PUJ} local access roads. The use of demand-oriented route construction is consistent with \cite{mandl1980}, whose urban transit network heuristic prioritizes the inclusion of high-demand nodes when constructing candidate public transport routes. In this study, the same principle is implemented computationally by drawing waypoints from the demand surface rather than from a uniform node distribution.

\textbf{Shortest-Path Connection and Loop Closure.}
After waypoints are selected, the generator connects consecutive waypoint pairs using \texttt{CityGraph.find\_shortest\_path()}. This method performs A$^\ast$ search over the CityGraph, but it evaluates only edges marked as \texttt{is\_drivable=True}. A$^\ast$ is suitable for this task because it combines accumulated path cost with a heuristic estimate of remaining distance, allowing efficient shortest-path computation while preserving optimality when the heuristic is admissible \citep{hart1968}. In the route-construction context, the resulting path between waypoint pairs is therefore the shortest feasible path through the \ac{PUJ}-eligible arterial subset.

The route is then closed into a loop by connecting the final waypoint back to the first waypoint. This loop structure is required because a jeepney route in the simulation is represented as a continuously traversable sequence of directed edges. If a valid drivable path cannot be found between any required waypoint pair, the attempted route is rejected and the generator resamples waypoints up to a configurable retry limit. The use of shortest-path-based procedures as route-construction components is consistent with urban transportation network design literature, where shortest-path and path-building heuristics are commonly used to generate candidate network alternatives \cite{farahani2013review}.

\textbf{Layer 2 Promotion.}
Once a continuous loop has been established on the Layer~0 infrastructure graph, the selected base edges are promoted into a transit-route representation. Promotion duplicates each selected infrastructure edge into a new directed edge whose start and end nodes are assigned to Layer 2. These promoted edges are then linked circularly through their \texttt{next\_edges} pointers so that each edge points to the next edge in the route, and the final edge points back to the first. This produces the route object traversed by jeepney agents during simulation.

This promotion step is intentionally a copy operation rather than an in-place modification of the CityGraph. The Layer~0 infrastructure graph remains the stable physical reference, while the promoted Layer~2 edges represent a candidate \ac{PUJ} service loop derived from that infrastructure. This separation allows many candidate routes and route systems to be generated, modified, and evaluated without corrupting the underlying city road network.

\textbf{Validation.}
Every generated \texttt{Route} object is validated before it is accepted. Three checks are required.

\begin{itemize} [nosep]
    \item First, \textbf{loop continuity} requires that the end node of each directed edge matches the start node of the next directed edge, including the connection from the final edge back to the first edge. This prevents broken route geometries.
    \item Second, \textbf{layer isolation} requires every edge in the accepted route to belong strictly to Layer 2 after promotion. This prevents a route object from accidentally mixing infrastructure edges with transit-route edges.
    \item Third, the \textbf{single out-pointer branching constraint} requires every route edge to have exactly one outgoing Layer 2 continuation. If an edge has more than one Layer 2 outgoing pointer, a jeepney agent would encounter an ambiguous branching decision while traversing what should be a deterministic route loop. Such routes are rejected.
\end{itemize}

Routes failing any of these checks are discarded and regenerated. Figure~\ref{fig:flowchart} provides a flowchart summarizing the complete route construction and validation pipeline.

This pipeline generates structurally valid, demand-responsive, arterially constrained route candidates for the optimization framework. It also ensures that all initial route chromosomes passed to the metaheuristic are continuous loops rather than disconnected edge collections.
 
\begin{figure}[H]
    \centering
    \includegraphics[width=0.8\textwidth]{chap3/figures/flowchart.jpg}
    \caption{Route Construction and Validation Pipeline}
    \label{fig:flowchart}
\end{figure}

\subsection{Coordinate Snapping via CKD Tree}

Coordinate snapping is required because raw geographic coordinates rarely coincide exactly with the topological nodes stored in the graph. Passenger origins, passenger destinations, and manually encoded route coordinates may fall between intersections or slightly away from the nearest stored road node. Without a snapping procedure, these coordinates could not reliably be used as pathfinding origins, pathfinding destinations, or route waypoints.

To resolve this spatial misalignment computationally, the framework employs the Compressed K-Dimensional Tree (\texttt{cKDTree}) implementation provided by the SciPy scientific computing library \citep{virtanen2020scipy}. Originally formalized as a general space-partitioning data structure by Bentley (1975), a KD-Tree organizes points in a $k$-dimensional space, specifically the two-dimensional geographic plane of longitudes and latitudes, by recursively bisecting the search area along alternating axes. 

The primary justification for utilizing a \texttt{cKDTree} rather than a standard geographic search is algorithmic efficiency. A naive nearest-neighbor procedure requires calculating the Haversine distance from the query coordinate to every single node in the network to find the minimum value. This requires an $\mathcal{O}(N)$ time complexity per query. In a parallelized agent-based simulation evaluating thousands of dynamic passenger coordinates across multiple generations, executing linear searches creates a severe computational bottleneck. By structuring the network nodes into a \texttt{cKDTree}, the nearest-neighbor query is reduced to a binary search through spatial partitions, decreasing the time complexity to $\mathcal{O}(\log N)$. This logarithmic scaling guarantees that coordinate snapping remains near-instantaneous regardless of how dense the underlying \texttt{CityGraph} becomes.

The implementation uses \texttt{cKDTree}-based snapping in two contexts:

The first context is route ingestion through \texttt{route\_from\_coords()}. This function supports externally defined route traces represented as \ac{JSON} coordinate sequences. The raw coordinates are converted into query points, snapped to the nearest valid CityGraph nodes using a \texttt{cKDTree} built from the infrastructure node coordinate array, and deduplicated to remove repeated consecutive matches. After snapping, A$^\ast$ shortest-path searches connect the snapped waypoints into a continuous path. The resulting path is then promoted into Layer~2 using the same route-construction protocol described in Section \ref{subsec:3.2.3 route construction and validation pipeline}. This allows externally prepared route coordinates to be converted into valid Route objects without requiring the input coordinates to coincide exactly with graph nodes.

The second context is passenger endpoint binding through \texttt{TravelGraph.\_snap\_node()}. When a passenger origin or destination is submitted to the TravelGraph, the coordinate may not match an existing layer-specific node. To resolve this, the TravelGraph constructs two coordinate arrays during initialization: one for Layer~1 nodes and one for Layer~3 nodes. Each array is indexed using SciPy's \texttt{cKDTree}.

The \texttt{\_snap\_node()} method receives a target Node and a target layer. If the target layer is~1, the method queries the Layer~1 tree and returns the nearest Layer~1 node. If the target layer is~3, it queries the Layer~3 tree and returns the nearest Layer~3 node. These snapped nodes become the valid A$^\ast$ origin and destination used by \texttt{findShortestJourney()}.

The method also preserves the project-wide coordinate convention \texttt{(lon, lat)}, which is necessary because reversing coordinate order would distort distance computations and nearest-neighbor matching. In addition, \texttt{findShortestJourney()} uses an LRU cache so that repeated \ac{OD} queries do not require redundant pathfinding computations.

Together, route-ingestion snapping and passenger-endpoint snapping ensure that all external coordinates are projected onto valid graph nodes before pathfinding begins. This keeps route construction, Journey computation, and simulation evaluation internally consistent because all movement is evaluated on the same topological road network.

\section{Direct Demand Model}
\label{sec:3.3 ddm}
\noindent{This study uses a \ac{DDM} to generate realistic origin and destination points across the city graph. Instead of sampling nodes uniformly, the model gives higher probability to locations that show stronger traffic activity and stronger network importance. It does this by querying a limited set of points from the TomTom Traffic Flow \ac{API}, filling the missing values through spatial interpolation, combining traffic with betweenness centrality, and then converting the final scores into a fast sampling table. This approach supports large maps because the model does not need to query every node in the network, and it also keeps repeated sampling efficient during route generation.}

\subsection{Traffic Data Ingestion via TomTom API}
\noindent To gather empirical traffic data for the sampled centroids under discrete temporal scenarios, the framework utilizes a dedicated \texttt{TrafficClient} module that isolates external network requests from the core mathematical engine. To query predictive traffic fields for specific target dates and hours, the client utilizes TomTom’s routing engine via the \texttt{calculateRoute} endpoint. The utilization of web mapping application programming interfaces (APIs) is highly relevant for this study, as travel time, related distance, and routes can be acquired by background high-performance platforms yielding precise, accurate, and realistic geographical information \citep{munoz2021study}. Understanding and discovering these urban mobility patterns facilitates better urban planning and traffic resource management \citep{munoz2021study}.

Because a routing engine natively requires an explicit origin-destination coordinate pair rather than an isolated point location, a directional \textit{virtual probe segment} heuristic is applied. For each sampled infrastructure node, a synthetic destination coordinate is projected at a fixed spatial offset ($+0.001^{\circ}$ latitude). This methodology aligns with recent literature that retrieves travel distances and times between pairs of points at different time slots to estimate localized transport variations \citep{munoz2021study}. This application acts as a predictive mechanism evaluating the traffic congestion state using the Estimated Time of Arrival (ETA) across specific spatial segments \citep{zafar2020traffic}.

From a single API response generated along these virtual probe links, the system simultaneously extracts two critical metrics. Recent deployments of TomTom's routing architecture validate this exact pipeline, noting its capability in delivering per-segment benchmarks by natively computing and returning both real-time variables and historical reference baselines within the same response payload \citep{vladut2025urbanscore}. Specifically, the client captures the dynamic, traffic-aware travel time ($t_i^{\text{travel}}$) from the API's \texttt{travelTimeInSeconds} parameter, and the theoretical free-flow travel time ($t_i^{\text{free}}$) from the \texttt{noTrafficTravelTimeInSeconds} parameter, as strictly defined within the response schema of the official TomTom Routing API \citep{tomtom2026routing}. In heterogeneous traffic models, the free-flow state is formally characterized as the desired speed that drivers can drive without being obstructed or influenced by other road users \citep{leong2020development}. It represents the travel time of a vehicle when its movement is not interfered with by other vehicles or interrupted by control devices \citep{leong2020development}. 

To quantify local congestion intensity, the model calculates the empirical traffic weight ($V_i$) as follows:

\[
V_i = \frac{t_i^{\text{travel}}}{t_i^{\text{free}}}
\]

This ratio perfectly mirrors established ETA-based congestion forecasting models, where the Congestion Index ($CI$) is calculated based on the current time for the road segment divided by the least (free-flow) time for the road segment \citep{zafar2020traffic}. Because different segments of the road network have varying physical properties, they must be evaluated based on their own isolated Congestion Index calculations \citep{zafar2020traffic}. A baseline weight of $V_i = 1.0$ reflects perfect free-flowing movement, whereas higher values strictly quantify the structural delay multiplier imposed by local traffic friction. This derived congestion index scales from zero up to discrete progressive labels corresponding to distinct traffic states: smooth, slightly congested, congested, highly congested, or blockage \citep{zafar2020traffic}. 

To eliminate redundant network operations and mitigate API throttling during iterative simulation loops, all network payloads are persistently cached. The cache keys are generated as an \ac{MD5} hash of a combined input string consisting of the node's geographic coordinates and its assigned ISO-8601 target datetime string, preserving the exact temporal snapshot parameters across identical simulation cycles.

\subsection{Spatial Interpolation via IDW and Haversine}
\noindent{After the model collects traffic weights for the sampled points, it estimates the values of the unqueried nodes using Inverse Distance Weighting (IDW). The model gives more influence to nearby sampled points and less influence to distant ones:

\[
W_j =
\frac{
\sum_{i=1}^{n} \left( \frac{V_i}{d_{ij}^{p}} \right)
}{
\sum_{i=1}^{n} \left( \frac{1}{d_{ij}^{p}} \right)
}
\]

where \(W_j\) is the imputed traffic weight for the target node \(j\), \(V_i\) is the known weight at sampled node \(i\), \(d_{ij}\) is the distance between the target node and the sampled node, \(n\) is the number of sampled points, and \(p\) is the power parameter that controls how fast influence decreases with distance. A higher \(p\) gives more weight to nearby points. \citep{Snyder1987} supports this kind of spherical spatial calculation when the method must work on geographic coordinates rather than flat coordinates.

The model computes \(d_{ij}\) using the Haversine formula instead of flat Euclidean distance. This matters because the road network lies on the surface of the Earth, not on a flat plane. The Haversine formula uses latitude and longitude and gives the great-circle distance between two points:

\[
a = \sin^2\left(\frac{\Delta \phi}{2}\right)
+ \cos(\phi_1)\cos(\phi_2)
\sin^2\left(\frac{\Delta \lambda}{2}\right)
\]

\[
c = 2 \cdot \operatorname{atan2}\left(\sqrt{a}, \sqrt{1-a}\right)
\]

\[
d = R \cdot c
\]

where \(\phi_1\) and \(\phi_2\) are the latitudes of point 1 and point 2, \(\Delta \phi\) is the difference between the two latitudes, \(\Delta \lambda\) is the difference between the two longitudes, \(R \approx 6{,}371{,}000\) meters is the mean radius of the Earth, and \(d\) is the final computed distance. This method stays numerically stable for small distances and avoids the precision problems that come from using the spherical law of cosines on nearby points.}

\subsection{OD Centrality Fusion}
\noindent The study combines traffic information with structural network importance to produce the final direct demand probability. Freeman (1977) introduced betweenness centrality as a measure of node importance based on how frequently a node lies on the shortest paths connecting other pairs of nodes in a network. The measure was designed to identify nodes that occupy strategically important positions in supporting movement and connectivity across the network. In this study, the same idea applies to the road system, where intersections that lie on many shortest paths serve as important connectors between different areas of the city, making betweenness centrality an appropriate structural basis for demand estimation.

\[
C_B(v) =
\sum_{\substack{s \neq v \neq t \\ s,t \in V}}
\frac{\sigma_{st}(v)}{\sigma_{st}}
\]

where \(C_B(v)\) is the betweenness centrality of node \(v\), \(\sigma_{st}\) is the total number of shortest paths between nodes \(s\) and \(t\), and \(\sigma_{st}(v)\) is the number of those paths that pass through \(v\). This gives a stable measure of structural importance in the transport network because it depends on shortest-path connectivity rather than on temporary traffic conditions.

Following the concept of \ac{OD} centrality, the model combines the traffic term and the centrality term into a single node score, with traffic weights estimated from sampled TomTom data and interpolated values for unqueried nodes, and structural importance represented by betweenness centrality \citep{Lowry2014}.

\[
S_i = W_i^{\alpha} \times C_i^{\beta}
\]

where \(W_i\) is the traffic-based weight of node \(i\), and \(C_i\) is its structural weight. The exponent \(\alpha\) controls the influence of traffic intensity, while \(\beta\) controls the influence of network structure. A larger value of \(W_i\) or \(C_i\) increases the final score, while a lower value in either term reduces it. In this way, the score favors nodes that are both busy and structurally important.

The raw scores are then normalized into probabilities using

\[
P_i = \frac{S_i}{\sum_{k=1}^{N} S_k}
\]

where \(P_i\) is the final probability assigned to node \(i\), \(S_i\) is the unnormalized direct demand score, and \(N\) is the total number of candidate nodes. This normalization ensures that all node probabilities sum to 1, allowing the values to be used directly for stochastic sampling. The resulting probabilities are then inserted into the Walker's Alias table for efficient route generation.

\subsection{Walker’s Alias Table Construction}
\noindent{After the model computes the final node probabilities, it stores them in a Walker Alias table to enable fast repeated sampling. Walker’s alias method constructs two arrays: the probability table and the alias table, in \(O(N)\) preprocessing time and then allows each sample to be drawn in \(O(1)\) time. This makes it efficient for large-scale simulations where the sampler is called many times, such as in route generation and passenger synthesis. Using this structure prevents performance bottlenecks as the network size increases while keeping sampling consistently fast. \citep{Walker1977} introduced this algorithm as an efficient method for generating discrete random variables from arbitrary probability distributions, which forms the basis of this optimization.

The construction process follows these steps:

\[
\sum_{i=1}^{N} P_i = 1
\]

The model scales the probabilities so that their mean becomes 1.0. It then separates the nodes into underfull and overfull groups. It fills each underfull entry using one alias entry from an overfull node until every slot reaches the same normalized scale. After that, one random index and one random float are enough to sample a node in constant time.

This optimization matters because the passenger generator and route search modules call the demand sampler many times during a run. With the alias table, the system avoids repeated binary searches and keeps the simulation fast even when it processes many candidate routes.}

\subsection{DDMConfig and Model Parameters}
\noindent{The \ac{DDM} uses fixed parameters to control how traffic weights and network centrality combine to form the final demand surface. The main parameters are:

\[
\alpha, \beta, \text{idw\_power}, \text{cache\_dir}
\]

The parameter \(\alpha\) controls the exponent applied to traffic weights \(W_i\). Higher values amplify the influence of congestion on demand. The parameter \(\beta\) controls the exponent applied to centrality scores \(C_i\). Higher values amplify the influence of network structure on demand. Together, they control the combined score as shown in the \ac{OD} Centrality Fusion section:

\[
S_i = W_i^{\alpha} \times C_i^{\beta}
\]

The parameter \(\text{idw\_power}\) controls the power exponent in the Inverse Distance Weighting formula used for traffic imputation. Higher values create sharper distance decay; lower values produce smoother interpolation. The \(\text{cache\_dir}\) setting specifies where to store cached TomTom \ac{API} responses, keeping them organized and making repeated runs reproducible.

All parameters remain fixed throughout simulation runs. Changing them between route candidates would alter the underlying demand surface and confound the route optimization. After initial selection, the same \(\alpha\), \(\beta\), and \(\text{idw\_power}\) values are used across all evaluations to ensure stable demand and fair route comparison. The demand sampler serves as the synthesis engine for the passenger generator, so consistency of these parameters is critical for reliable simulation behavior.

\section{The Three-Layer Graph Architecture}
\label{sec:3.4 three layer graph}
The spatial network, route system, and calibrated behavioral weights are integrated through the \texttt{TravelGraph}. The \texttt{TravelGraph} is the central graph structure used to generate passenger \texttt{Journey} objects in the simulation. It is constructed from two inputs: the \texttt{CityGraph} discussed in Section \ref{sec:3.2 CityGraph}, which represents the physical road network of Iligan City, and a Route System, which contains the validated jeepney routes operating on that network. The \texttt{CityGraph} provides the underlying geographic coordinates, road nodes, and directed road segments, while the Route System supplies the transit paths that passengers may board and ride.

In this study, a \texttt{Journey} refers to the ordered sequence of directed edges followed by a passenger from origin to destination. This sequence may include walking, waiting, riding, alighting, transferring, and direct walking. The \texttt{TravelGraph} is therefore not only a map representation. It is a pathfinding-ready behavioral graph that converts the road network and jeepney routes into a structure where passenger movement can be evaluated using generalized travel cost.

The \texttt{TravelGraph} is organized into three passenger-state layers. Layer 1 represents the start-walk network, where passengers walk from their origin toward possible boarding points. Layer 2 represents the ride network, where passengers move only along jeepney routes. Layer 3 represents the end-walk network, where passengers walk after alighting and where transfers are represented. These layers are spatially co-registered, meaning that the same physical coordinate may appear in multiple layers, but each layer represents a different travel state.

This structure allows the shortest-path algorithm to evaluate complete multimodal passenger journeys without confusing walking movement with riding movement. A passenger cannot ride unless the path enters Layer 2 through a valid wait or transfer edge. A passenger cannot leave a vehicle unless the path exits Layer 2 through an alight edge. Thus, the graph structure itself enforces the behavioral rules of jeepney travel.

\subsection{Three-Layer Graph Architecture Rationale}
A single-layer graph is sufficient for ordinary road shortest-path problems, but it is not sufficient for modeling jeepney passenger movement. In a flat graph, a road intersection is represented by only one node. This means that a walking passenger and a riding passenger located at the same intersection would share the same outgoing edges, even though they are in different travel states. This would allow unrealistic behavior, such as walking directly from a ride path without alighting, or switching between jeepney routes without a transfer event.

The three-layer architecture resolves this by separating travel states into different topological layers. A passenger walking to a boarding point is represented on Layer 1. A passenger riding a jeepney is represented on Layer 2. A passenger walking after alighting or preparing to transfer is represented on Layer 3. Since each layer has its own node instances and edge types, the outgoing movements available to a passenger depend on the passenger's current state.

This design follows the logic of multilayer transport network modeling. Multimodal transport systems are better represented as interconnected layers rather than as a single aggregated graph because different modes and operational states have different constraints, costs, and movement rules \citep{alessandretti2023multimodal}. \cite{peng2023multi} similarly demonstrate that multilayer transport representations preserve distinctions between different transport flows that would be obscured in an aggregated network. In this study, the same principle is applied to jeepney passenger routing: walking, riding, alighting, and transferring are not treated as identical movements, but as distinct behavioral states encoded directly into the graph.

The three-layer structure also preserves the correctness of standard shortest-path search. Since waiting, riding, walking, alighting, transferring, and direct walking are all represented as directed edges with explicit weights, the search algorithm does not need special route-choice rules outside the graph. The \texttt{TravelGraph} itself ensures that every computed \texttt{Journey} is physically valid and behaviorally interpretable.

\begin{figure}[H]
  \centering
  \includegraphics[width=0.8\textwidth]{chap3/figures/travel_graph_visualization.jpg}
  \caption{Three-Layer Graph Architecture used for constructing TravelGraph objects}
  \label{fig:iligan_network}
\end{figure}

\subsection{Layer Initialization and Structure}
\noindent
The \texttt{TravelGraph} initializes three layers from the \texttt{CityGraph} and the Route System. Each layer uses the same geographic coordinate system, following the project-wide coordinate convention of longitude and latitude. This consistency is necessary because all edge lengths are computed geographically and all snapping procedures depend on correctly ordered coordinates.

The edge-weight equations in this subsection use the calibrated \texttt{TravelGraph} cost weights reported in Table~\ref{tab:final_calibrated_cost_weights}. 

\begin{itemize}
    \item \textbf{Layer 1 is the start-walk layer.} For every node in the \texttt{CityGraph}, the \texttt{TravelGraph} creates a corresponding Layer 1 node at the same coordinate. The road edges from the \texttt{CityGraph} are then copied into Layer 1 as start-walk edges, denoted as \texttt{SW} edges. These edges represent the initial walking movement from a passenger's origin toward possible boarding points. The cost of each start-walk edge is computed as:

    \begin{equation}
        w_{\mathrm{SW}} = d \cdot \text{\texttt{walk\_wt}}
    \end{equation}

    where \(d\) is the edge length in meters and \(\text{\texttt{walk\_wt}}\) is the calibrated walking cost in \ac{EIVM} per meter.

    \item \textbf{Layer 2 is the ride layer.} Unlike Layer 1 and Layer 3, Layer 2 is not a full duplicate of the road network. It contains only the route edges belonging to validated jeepney routes. Each route is represented as a separate cyclic subgraph composed of ride edges, denoted as \texttt{RI} edges. The cost of each ride edge is computed as:

    \begin{equation}
    w_{\mathrm{RI}} = d \cdot \text{\texttt{ride\_wt}}
    \end{equation}

    where \(d\) is the route edge length in meters and \(\text{\texttt{ride\_wt}}\) is the calibrated ride cost in \ac{EIVM} per meter.

    \item \textbf{Layer 3 is the end-walk layer.} For every node in the \texttt{CityGraph}, the \texttt{TravelGraph} creates a corresponding Layer 3 node at the same coordinate. The \texttt{CityGraph} road edges are copied into Layer 3 as end-walk edges, denoted as \texttt{EW} edges. These edges represent the walking movement after alighting from a jeepney or after entering the destination-side pedestrian layer. The cost of each end-walk edge follows the same walking cost function:

    \begin{equation}
    w_{\mathrm{EW}} = d \cdot \text{\texttt{walk\_wt}}
    \end{equation}

    where \(d\) is the edge length in meters and \(\text{\texttt{walk\_wt}}\) is the same calibrated walking cost used for Layer 1.
\end{itemize}

Together, the start-walk, ride, and end-walk layers allow the model to represent a complete passenger journey as a sequence of state-specific movements. The passenger begins in Layer 1, may enter Layer 2 by boarding, may exit to Layer 3 by alighting, and may continue walking or transfer through Layer 3 depending on the computed least-cost path.

\subsection{Inter-Layer Transition Edges}
The three layers are connected through four directed inter-layer transition edges: wait, alight, transfer, and direct. These transition edges do not represent road segments. Instead, they represent behavioral events that move a passenger from one travel state to another.

\begin{itemize}
    \item \textbf{Wait edges.} Wait edges, denoted as \texttt{WA} edges, connect Layer 1 to Layer 2. A wait edge allows a passenger to board a jeepney route from the start-walk layer. Its weight is:

    \begin{equation}
    w_{\mathrm{WA}} =
    \text{\texttt{wait\_wt}}
    \end{equation}

    where \text{\texttt{wait\_wt}} is the calibrated waiting or initial boarding cost. This edge represents the burden of waiting for and boarding a jeepney at the first boarding point.

    \item \textbf{Transfer edges.} Transfer edges, denoted as \texttt{TR} edges, connect Layer 3 back to Layer 2. A transfer edge allows a passenger who has already alighted to board another jeepney route. Its weight is:

    \begin{equation}
    w_{\mathrm{TR}} =
    \text{\texttt{transfer\_wt}}
    \end{equation}

    where \text{\texttt{transfer\_wt}} is the calibrated transfer penalty. This weight captures the additional disutility of changing jeepneys beyond walking and riding cost. The rationale for the numerical value of \text{\texttt{transfer\_wt}} is discussed in Section \ref{subsec:3.1.4 parameter extraction and derivation}.

    \item \textbf{Alight edges.} Alight edges, denoted as \texttt{AL} edges, connect Layer 2 to Layer 3. An alight edge allows a passenger to leave a jeepney and enter the end-walk layer. Its weight is:

    \begin{equation}
    w_{\mathrm{AL}} = 0
    \end{equation}

    This does not mean that alighting is ignored. Rather, the alight edge is a structural state-change connector. The ride cost has already been charged through ride edges, and any walking after alighting is charged through end-walk edges. If a passenger transfers after alighting, the transfer penalty is charged separately through the transfer edge. Assigning a positive cost to alighting would therefore double-count travel burden already represented by other terms.

    \item \textbf{Direct edges.} Direct edges, denoted as \texttt{DI} edges, connect Layer 1 to Layer 3 at the same coordinate. A direct edge allows the model to evaluate a pure walking alternative without requiring jeepney use. Its weight is:

    \begin{equation}
    w_{\mathrm{DI}} = 0
    \end{equation}

    This edge is also structural. It does not represent free travel. Walking cost is already charged through \texttt{SW} and \texttt{EW} walking edges. The direct edge simply permits the pathfinding algorithm to move from the origin-side walking layer to the destination-side walking layer when a non-transit path is cheaper than a jeepney-assisted path.
\end{itemize}

Thus, \texttt{alight\_wt} and \texttt{direct\_wt} are both set to zero because they are not independent behavioral penalties. They are state-transition connectors that prevent the graph from forcing artificial costs onto movements whose burdens have already been accounted for through walking, riding, waiting, or transfer terms. This follows generalized cost decomposition, where each component of travel burden should be counted once through its proper term \cite{ortuzar2011modelling}. It is also consistent with transit assignment models that represent access, ride, alight, and transfer links separately, with cost assigned only when the link represents an independent source of impedance \cite{spiess1989optimal}.

\subsection{Route Isolation and Graph Stitching}
After the layers and transition edges are created, the \texttt{TravelGraph} must connect valid edge sequences so that pathfinding can proceed across the network. This is done through graph stitching. The stitching procedure links one directed edge to another when the end node of the first edge matches the start node of the second edge by coordinate and layer. In the implementation, the \texttt{\_stitch()} function uses a hash-based coordinate lookup from \((\texttt{lon}, \texttt{lat}, \texttt{layer})\) to outgoing edges. This avoids repeatedly scanning all possible edge pairs and allows adjacency construction through constant-time lookup for matching coordinates.

Figure~\ref{fig:CityGraph_to_TravelGraph} summarizes the full construction sequence from the CityGraph and validated Route list to the final pathfinding-ready TravelGraph. The final stages of the pipeline show the graph-stitching operations that connect valid edge sequences while preserving route-specific Layer 2 isolation.

\begin{figure}[H]
  \centering
  \includegraphics[width=0.8\textwidth]{chap3/figures/city_to_travel_flowchart.jpg}
  \caption{CityGraph to TravelGraph construction and route-isolated stitching pipeline}
  \label{fig:CityGraph_to_TravelGraph}
\end{figure}

However, stitching must be handled carefully for Layer 2. Since multiple jeepney routes may pass through the same physical intersection, a naive stitching procedure that connects all Layer 2 ride edges at the same coordinate would create a teleportation anomaly. In that anomaly, a passenger riding Route A could instantly continue on Route B at a shared intersection without alighting, transferring, waiting, or paying the proper behavioral cost. This would make the computed \texttt{Journey} physically invalid and would underestimate the cost of transfers.

To prevent this, Layer 2 edges are scoped strictly by route index. The \texttt{TravelGraph} creates separate Layer 2 node populations for each route, even when two routes share the same geographic coordinate. During stitching, ride edges are connected only to ride edges belonging to the same route index. Wait and alight edges are also stitched only to the Layer 2 nodes associated with the specific route they connect to. This preserves route containment: once a passenger enters a route's ride layer, the passenger remains on that route until an explicit alight edge moves the path into Layer 3.

If the passenger later boards another route, the path must traverse a transfer edge from Layer 3 back to Layer 2. This guarantees that route changes are counted as transfer events and that \text{\texttt{transfer\_wt}} is applied exactly when required. The result is a graph where shared physical intersections do not automatically imply shared vehicle states.

This route isolation is essential for jeepney network simulation because jeepney routes frequently overlap along major corridors. Without route-index isolation, overlapping corridors would incorrectly allow passengers to switch vehicles at no cost. With isolation, overlapping routes remain geographically co-located but operationally distinct, which better reflects real passenger behavior and the physical constraint that a commuter cannot change jeepneys without first alighting and boarding again.

\subsection{A* Pathfinding and Admissibility Proof}
\label{subsec:3.4.5 astar_proof}
\noindent
To simulate realistic passenger movement across the \texttt{TravelGraph}, the framework must continuously compute the minimum-cost multimodal path between dynamic origin and destination coordinates. Because the graph evaluates thousands of agent decisions per simulation tick, executing Dijkstra's algorithm, which explores the graph uniformly in all directions, would create a severe computational bottleneck. The framework therefore utilizes the A* search algorithm, which accelerates pathfinding by employing a heuristic function to guide the search directionally toward the target \citep{russell2010artificial}.

For A* search to mathematically guarantee the discovery of the true optimal path, the heuristic function $h(n)$ must be strictly admissible \citep{russell2010artificial}. An admissible heuristic is one that never overestimates the true remaining cost to reach the destination from any given node $n$. Formally, the algorithm evaluates nodes based on the function:

$$f(n) = g(n) + h(n)$$

where $g(n)$ is the exact accumulated cost from the start node to node $n$, and $h(n)$ is the estimated remaining cost from node $n$ to the goal node $t$. Let $h^*(n)$ denote the true minimal cost to reach $t$ from $n$. The admissibility condition requires that:

$$\forall n, \quad h(n) \leq h^*(n)$$

In the \texttt{TravelGraph}, the true cost $h^*(n)$ of any path $P$ is the sum of the generalized costs of its constituent edges. The path $P$ can be decomposed into spatial edges $P_{\mathrm{spatial}}$ (start-walk, ride, and end-walk edges) that cover physical geographic distance, and structural transition edges $P_{\mathrm{events}}$ (wait, transfer, alight, direct) that represent state changes without geographic displacement:

$$h^*(n) = \sum_{e \in P_{\mathrm{spatial}}} (d_e \cdot \beta_e) + \sum_{k \in P_{\mathrm{events}}} W_k$$

where $d_e$ is the geographic length of spatial edge $e$, $\beta_e$ is the cost-per-meter parameter (either $\beta_{\mathrm{walk}}$ or $\beta_{\mathrm{ride}}$), and $W_k$ represents the fixed penalty for behavioral events (where $W_k \geq 0$). 

To construct an admissible heuristic, we must establish the absolute lower bound of this cost. Since all structural transition penalties are non-negative ($W_k \ge 0$), discarding them provides a strict inequality:

$$h^*(n) \geq \sum_{e \in P_{\mathrm{spatial}}} (d_e \cdot \beta_e)$$

Furthermore, for any spatial edge $e$, the cost parameter $\beta_e$ must be greater than or equal to the minimum traversal weight available in the entire network. Let $\beta_{\mathrm{min}} = \min(\beta_{\mathrm{walk}}, \beta_{\mathrm{ride}})$. By substituting $\beta_{\mathrm{min}}$, the inequality is preserved:

$$h^*(n) \geq \sum_{e \in P_{\mathrm{spatial}}} (d_e \cdot \beta_{\mathrm{min}})$$
$$h^*(n) \geq \beta_{\mathrm{min}} \sum_{e \in P_{\mathrm{spatial}}} d_e$$

The sum $\sum d_e$ represents the total geographic distance traversed along the road and transit network. According to the triangle inequality in metric spaces, the sum of the segments of any network path between two points must be greater than or equal to the direct straight-line distance between those points. Let $D_{\mathrm{straight}}(n, t)$ be the Haversine great-circle distance between node $n$ and the destination $t$. Therefore:

$$\sum_{e \in P_{\mathrm{spatial}}} d_e \geq D_{\mathrm{straight}}(n, t)$$

Substituting this geometric constraint back into the cost equation yields the absolute lower bound of the remaining path cost:

$$h^*(n) \geq D_{\mathrm{straight}}(n, t) \cdot \min(\beta_{\mathrm{walk}}, \beta_{\mathrm{ride}})$$

This lower bound is exactly the formulation utilized by the framework's A* heuristic:

$$h(n) = D_{\mathrm{straight}}(n, t) \cdot \min(\beta_{\mathrm{walk}}, \beta_{\mathrm{ride}})$$

Because the straight-line geographic distance is geometrically minimal, and the selected traversal weight represents the absolute minimum cost-per-meter available in the environment, the product $h(n)$ can never exceed the true cost $h^*(n)$ that must be routed through the constrained road network. 

Since $h(n) \leq h^*(n)$ holds unconditionally for all nodes in the \texttt{TravelGraph}, the heuristic is strictly admissible. This structural proof guarantees that the A* algorithm will always evaluate the true optimal multimodal journey for every simulated passenger without compromising computational efficiency.

\subsection{Differentiated Wait and Transfer Weights Rationale}
The \texttt{TravelGraph} distinguishes between initial waiting and route transfer because these events represent different forms of passenger burden. Initial waiting occurs when a passenger boards the first jeepney route from the start-walk layer. A transfer occurs only after a passenger has already ridden one route, alighted into the end-walk layer, and boarded another route. Although both events involve entering a jeepney, their behavioral meaning is not identical.

The wait weight represents the generalized burden of initial boarding and expected waiting. The transfer weight represents the additional disutility of interrupting a trip, navigating to another boarding opportunity, facing uncertainty about the next vehicle, and losing the continuity of a direct ride. Prior public transport studies show that transfer penalties remain substantial even after walking and waiting components are controlled. \cite{jara-diaz2022} propose a planning range of 13 to 18 \ac{EIVM} for the pure transfer penalty, while \cite{garciamartinez2018transfer} estimate pure transfer penalties in the range of approximately 15.2 to 17.7 \ac{EIVM}. \cite{yap2024} further show that passengers value bus-bus interchanges more negatively than metro-metro interchanges, partly because of comfort, reliability, frequency, and perceived safety factors. \cite{meskell2024} similarly emphasize that bus-bus transfers are affected by uncertainty, exposure to weather, road crossing, and service reliability.

For this reason, the \texttt{TravelGraph} applies \text{\texttt{wait\_wt}} only to \texttt{WA} edges and \text{\texttt{transfer\_wt}} only to \texttt{TR} edges. This prevents the model from treating first boarding and subsequent transfers as identical events. It also prevents double-counting, since alight edges carry zero cost and transfer penalty is charged only when the passenger boards another route after alighting. The structure therefore keeps the generalized cost model behaviorally realistic while preserving compatibility with standard shortest-path search.

\section{Agent-Based Passenger Simulation}
\noindent
The agent-based passenger simulation models commuters as autonomous agents that generate trips, select feasible journeys, and move through the multi-layer travel graph under explicit behavioral and operational rules. The simulation serves as the runtime core of the framework because it links the demand model, the travel graph, and the jeepney fleet into a single dynamic commuting environment.

\subsection{Passenger Generation via the Direct Demand Model}
\noindent
Passenger generation begins with the \ac{DDM}, which estimates where trips are likely to originate and end across the study area. This architecture significantly upgrades the baseline framework established by \citep{sanchez2025}, which relied on a simplified graph network with uniform passenger distribution and a simple division of zones into Residential and Non-Residential areas. In contrast, this model establishes spatial heterogeneity by combining \ac{GIS} infrastructure with real-world traffic flows to calculate \ac{OD} centrality.

The model assigns a probability to each valid node by combining traffic-based spatial weight and network structure. For a node $i$, the final demand score is expressed as:
\[
P_i = W_i^\alpha \times C_i^\beta
\]
where $W_i$ is the empirical or interpolated traffic weight, $C_i$ is the betweenness centrality score derived from \citep{Freeman1977}, and $\alpha$ and $\beta$ are calibration parameters that control the relative influence of live traffic and network topology. This formulation adapts \citep{Lowry2014} \ac{OD} centrality framework by substituting static census data with dynamic \ac{API} data. Because resource constraints require a random sampling procedure of the local road network \citep{BowlingAultmanHall2003}, the final resolved probabilities are structured into an alias table \citep{Walker1977} to guarantee \(O(1)\) time complexity during subsequent agent generation.

The traffic weight is derived from the ratio of free-flow speed to current traffic speed:
\[
W_i = \frac{v_{\text{free}}}{v_{\text{current}}}
\]
where $v_{\text{free}}$ is the free-flow speed and $v_{\text{current}}$ is the observed or estimated traffic speed. The structural importance of each node $i$ is computed using betweenness centrality \citep{Freeman1977}:
\[
C_i = \sum_{s \neq i \neq t \in V} \frac{\sigma_{st}(i)}{\sigma_{st}}
\]
where $\sigma_{st}$ is the total number of shortest paths between nodes $s$ and $t$, and $\sigma_{st}(i)$ is the number of those paths that pass directly through node $i$.

The resolved probabilities are stored in Walker's Alias Table so that spatial sampling can run in constant time during simulation \citep{Walker1977}. This makes repeated demand generation efficient even when the optimization module evaluates many candidate route systems.

Once spatial passenger demand is calculated via the \ac{DDM}, the programmatic distribution of these generated passenger journeys is handled deterministically via the schedule generation routine of the passenger generator module. To avoid structural spawning bottlenecks on the first simulation tick, the schedule smoothly distributes the sampled \ac{OD} pairs across a discrete 100-tick window using a Gaussian standard deviation. Manual testing across discrete passenger volume levels demonstrated high consistency and framework stability. This consistency justifies the omission of highly complex, time-varying arrival rates or temporally inhomogeneous Poisson processes, allowing the model to remain structurally streamlined while strictly retaining its essential spatial heterogeneity.

After the origin and destination are sampled, the framework computes the passenger journey on the TravelGraph using the A* shortest-path algorithm proven admissible in Section \ref{subsec:3.4.5 astar_proof}.

\subsection{Passenger State Machine}
\noindent
Passenger behavior follows a \ac{FSM} restricted to exactly four valid stages: walking, waiting, riding, and done. The state transition mechanism controls these transitions by evaluating the passenger's current behavioral state at each discrete simulation tick and executing the corresponding geographic or topological updates.

\begin{figure}[h]
  \centering
  \includegraphics[width=0.8\textwidth]{chap3/figures/passenger_graph.jpg}
  \caption{Passenger Graph Traversal Visualization}
  \label{fig:passenger_graph_traversal}
\end{figure}

\textbf{Walking State.} In the walking state, the passenger behaves as an active pedestrian traversing the multi-modal network along the Layer 1 (Start Walk), Layer 3 (End Walk), or Direct ($DI$) transition edges. Movement during a single simulation tick is determined by the kinematic equation:
\begin{equation}
d = v_{\text{walk}} \times \Delta t
\end{equation}
where $v_{\text{walk}}$ represents the scalar walking velocity and $\Delta t$ is the discrete simulation time delta. Here, it must be re-iterated that Direct transition edges are implemented as length 0 and are therefore ineffectual.

The spatial progression tracking routine tracks progression linearly in meters along the passenger's current graph edge. If the calculated distance $d$ traveled within a single tick exceeds the remaining length of the current edge, the execution sequence executes a sub-tick edge wrapping sequence. The system subtracts the remaining edge length from $d$, increments the journey path index to transition the passenger to the next contiguous edge within that same tick, and immediately applies the leftover distance to the new edge.

When an edge transition shifts the passenger onto a Wait ($WA$) transition edge, the system halts walking kinematics and automatically triggers a state change to the waiting state.

\textbf{Waiting State.} Upon entering the waiting state, the passenger becomes completely idle at a designated transit stop node on Layer 1. During this phase, the passenger's local progression mechanics are paused, and the agent accumulates wait time penalties within the environment. The agent remains locked in this idle state until the master simulation controller emits a boarding event, which occurs when a Jeep agent with available physical capacity arrives at the matching topological node. The simulation controller then transfers the passenger into the vehicle's internal occupancy registry and updates the passenger's state directly to riding.

\textbf{Riding State.} In the riding state, the passenger is situated inside a Jeep agent operating exclusively on Layer 2 of the multi-layer travel graph. The passenger's physical movement is completely bound to the vehicle, meaning local kinematic updates are paused for the agent. Instead, the passenger passively monitors the 1D topological progress of the host vehicle along its fixed transit loop. When the Jeep emits an edge transition event indicating it has reached the node that matches the passenger's pre-computed alighting node, the simulation controller handles the disembarkation logic, moves the passenger down to the pedestrian network, and reverts the state machine back to the walking state.

\textbf{Done State.} The done state represents the terminal absorption state of the passenger's goal-oriented path. The agent transitions to this state automatically upon completing the final pedestrian edge of their multi-modal journey sequence on the graph. Once flagged as arrived, the passenger agent is removed from active temporal processing loops, allowing the simulation upon completion to finalize and record the aggregate journey cost metrics.

\subsection{PUJ Agents and Systems}
\noindent
\ac{PUJ}s are modeled as autonomous vehicle agents that move continuously along fixed route loops. Each jeepney maintains its route, current edge index, progress along the edge, capacity, speed, and onboard passengers. The vehicle does not store a full geographic path at every tick. Instead, it uses topological progression along the route, which keeps the update process efficient.

During each simulation tick, the jeepney evaluates its spatial progress using the kinematic relation

\[
d = v \times dt
\]

where $v$ is the jeepney speed and $dt$ is the time step. If the accumulated progress exceeds the current edge length, the jeepney advances to the next edge in the loop. When the route reaches the end, the index wraps around to the beginning, which matches the cyclic nature of jeepney operations.

Boarding is allowed only when the jeepney still has free capacity:

\[
n_{\text{passengers}} < C_{\text{max}}
\]

where $n_{\text{passengers}}$ is the current occupancy and $C_{\text{max}}$ is the maximum vehicle capacity. The jeepney system groups all vehicles by route and updates them in batch. This keeps fleet management modular and prevents the simulation loop from handling each vehicle separately.

\subsection{Spatial Headway Distribution and Fleet Allocation via Mohring Effect}
\noindent
The framework avoids bus bunching by distributing jeepneys evenly along each route at initialization \citep{yang2020}. Bus bunching occurs when vehicles that initially operate with balanced spacing gradually cluster together because of unequal passenger loading and delay accumulation.

To reduce this issue, the framework spatially distributes jeepneys across the route loop before the simulation begins. If all vehicles start at the same point, they cluster quickly and create unstable headways. To prevent that, the initial spacing distance is computed as

\[
S = \frac{L_{\text{loop}}}{N}
\]

where $L_{\text{loop}}$ is the total length of the closed route loop and $N$ is the number of jeepneys assigned to that route. This spacing is measured equidistantly along the perimeter of the closed loop, so each jeepney is placed at a fixed progress interval on the same continuous route. \citep{mohring1972}

Fleet size is dynamically allocated across the network using the Mohring effect, which states that transit service frequency should scale with ridership to optimize overall passenger waiting times \citep{mohring1972}. In this framework, the \textit{demand level} ($D_r$) of a route refers specifically to the aggregate passenger-distance generated by sampled passenger agents whose shortest-path journeys (determined via the multi-layer TravelGraph) actively utilize that specific route's transit edges. Rather than simply counting the raw volume of independent agents, this edge-weighted approach ensures that longer, high-capacity commuter corridors receive proportionally more fleet priority than short shuttle routes with equivalent boarding counts. To translate this demand level into discrete fleet counts, a minimum demand threshold is enforced to prevent complete route starvation before a square-root allocation rule is implemented:

\[
f_r = F_{\text{tot}} \cdot \frac{\sqrt{D_r}}{\sum_{k=1}^{m}\sqrt{D_k}}
\]

where $f_r$ is the allocated fleet share of route $r$, $F_{\text{tot}}$ is the total available fleet across the city network, and $D_r$ is the localized demand level of that route. This formulation keeps service frequency mathematically responsive to simulated commuter density while preserving network-wide resource balance \citep{mohring1972}.

By combining spatial headway balancing with demand-responsive fleet allocation, the framework reduces clustering behavior while improving service coverage across the simulated transit network.

It is important to distinguish this shortest-path sampling methodology from a simple linear summation of the \ac{DDM}. While directly summing the geographic demand probabilities of the nodes along a route is computationally trivial, it only measures \textit{local geographic density}---essentially, how busy the streets are that a route traverses. This static approach fails to account for topological connectivity, potentially overallocating vehicles to poorly designed routes that pass through dense areas but fail to connect logical origins to destinations. Within the formal \ac{TRNDP}, true network utility cannot be derived from isolated spatial coverage; it must be evaluated by how efficiently the route system satisfies many-to-many travel demand across continuous \ac{OD} paths \citep{kepaptsoglou2009transit}. By evaluating demand via simulated shortest-path journeys across the multi-layer \texttt{TravelGraph}, the framework captures true \textit{network-level utilization}. This dynamically accounts for \ac{OD} matching, transfer synergies, and multi-modal network flow, which are critical requirements for accurately modeling public transit passenger assignment \citep{farahani2013review}. Consequently, this ensures that fleet capacity is strictly allocated based on realized commuter utility rather than mere geographic proximity.

\subsection{Simulation and Event Architecture}
\noindent
The simulation module serves as the master controller of the framework. It initializes the TravelGraph, activates the PassengerGenerator, manages the JeepSystem, and executes the main time loop. Each simulation run begins from a clean setup so that no state leaks from one candidate route system to another.

At each tick, the simulation follows this order:
\begin{enumerate}
  \setlength{\itemsep}{0pt}
  \setlength{\parskip}{0pt}
  \setlength{\parsep}{0pt}
  \item Generate new passengers according to the demand process.
  \item Update all jeepney agents.
  \item Update all passenger agents.
  \item Process boarding and alighting events.
  \item Record the current metrics.
\end{enumerate}

This order matters because vehicle movement must be updated before passenger interaction. A passenger boards only after the jeepney has reached the relevant node for that tick. The simulation does not allow direct interaction between passengers and vehicles. Instead, jeepneys emit events, and the main controller handles boarding, capacity checking, and alighting through the TravelGraph. This event-driven structure keeps the model modular and computationally efficient.

\subsection{Simulation Fitness Evaluation}
\noindent
The optimization framework utilizes a single, authoritative fitness function to evaluate the actual dynamic performance of a route system under complete agent-based interaction. This simulation fitness is the exclusive metric used by the \ac{GA} for tournament selection, crossover ranking, and elitism. Unlike simplistic topological models, this metric captures the compounding effects of vehicle capacity limits and transfer delays. Following the established consensus in \ac{TRNDP} literature, the core objective function is the minimization of Total User Cost rather than average travel time \citep{kepaptsoglou2009transit, fan2006optimal}. Minimizing average time is structurally flawed for evolutionary algorithms, as it allows the optimizer to artificially improve its score by stranding difficult-to-serve passengers.

The comprehensive fitness function evaluates three distinct terms:

\[
F_{\text{sim}} = \sum_{i \in \mathcal{C}} T_i + \sum_{j \in \mathcal{I}} \left( T_j^{\text{elapsed}} + \beta \hat{T}_j^{\text{rem}} \right) + \alpha \sigma(T_i \mid i \in \mathcal{C})
\]

The first term is the \textbf{Total User Cost}, which sums the realized door-to-door generalized travel time ($T_i$) for all completed passengers ($\mathcal{C}$). Because the multi-layer \texttt{TravelGraph} already converts physical distance and transfer penalties into \ac{EIVM}, this sum inherently captures the full behavioral friction of the network \citep{ortuzar2011modelling}.

The second term is the \textbf{Underservice Penalty}, acting as an exterior penalty function to handle the implicit constraint that all simulated passengers must complete their journeys \citep{coello2002theoretical}. For incomplete passengers ($\mathcal{I}$), the algorithm adds their elapsed time ($T_j^{\text{elapsed}}$) to a multiplier ($\beta$) of their estimated remaining A* cost ($\hat{T}_j^{\text{rem}}$). With $\beta > 1$, an incomplete journey is strictly more expensive than a completed one, forcing the optimizer to prioritize spatial coverage.

The third term is the \textbf{Equity Regularizer}. Utilizing the population standard deviation ($\sigma$) of completed travel times scaled by a weight ($\alpha$), this term penalizes route systems that deliver highly unequal service distributions. This operationalizes horizontal transit equity, ensuring the optimizer does not sacrifice peripheral communities to marginally improve high-density trunk routes \citep{welch2013measure}.

\subsection{Static Surrogate Cost Approximation}\noindentWhile the simulation fitness provides a high-fidelity evaluation, executing the temporal physics engine requires updating thousands of independent agent states across hundreds of simulation ticks. During the Lamarckian local search phase (detailed in Section 3.6.2), the algorithm must evaluate multiple candidate route perturbations per chromosome. Running the full simulation for each minor structural mutation would severely bottleneck the evolutionary search.To resolve this computational limitation, the framework utilizes a static surrogate cost approximation. It is critical to note that this surrogate is not a fitness function; it is a lightweight heuristic evaluated strictly to bypass expensive objective calculations during intra-generational local searches \citep{jin2005survey}. The surrogate bypasses the temporal simulation entirely, directly computing the static, multi-modal A* generalized cost for a representative sample of \ac{OD} pairs ($M$), combined with a penalty for excessive route length ($L_{\text{route}}$):$$C_{\text{sur}} = \sum_{p=1}^{M} c^{\text{A*}}_p + \eta L_{\text{route}}$$To compensate for the absence of dynamic boarding delays and vehicle capacity constraints in the static approximation, the surrogate evaluation artificially inflates the topological transfer penalty by a factor of 2.5. This heuristic adjustment integrates empirical transfer aversion directly into the $\mathcal{O}(1)$ evaluation layer, heavily penalizing highly fragmented route configurations that may appear efficient statically but would fail under temporal simulation \citep{iseki2009transfers}.This surrogate cost is strictly confined to localized mutation operations. It functions exclusively as an $\mathcal{O}(1)$ acceptance gate: if a Lamarckian mutation structurally improves the surrogate cost, the route modification is provisionally kept. Once the local search phase concludes, the fully mutated chromosome is routed back through the microscopic temporal simulation to receive its authoritative $F_{\text{sim}}$ fitness score before the \ac{GA} advances to the next generation.

\section{The Hybrid GA-ACO Memetic Algorithm}
\noindent
The optimization layer combines a \ac{GA} with an Ant Colony
Optimization (ACO)-style memory structure to produce a memetic search procedure. In this framework, a chromosome represents an entire jeepney route system, while a gene represents one jeepney route inside that system. This distinction matters because the optimizer must preserve the integrity of closed transit loops, not merely rearrange isolated road segments. The \ac{GA} supplies global exploration through recombination and selection, while the ACO component supplies stigmergic
memory through pheromone trails that record where passenger demand and route utility are strongest. The resulting procedure is memetic because candidate solutions are not only inherited, but also locally improved before they are passed to the next generation \citep{moscato1989evolution, moscato2003memetic}.

The memetic operators execute within each generation of the search loop.
Specifically, after crossover produces a child chromosome and that child is evaluated for the first time, the Lamarckian local search operators are applied to the child before it enters selection for the next generation. This ordering is detailed in Section~\ref{sec:pipeline}, but the principle is worth establishing here: local improvement happens inside the generation, not between
generations.

% ────────────────────────────────────────────────────────────
\subsection{Stigmergic Pheromone Matrix and the Demand-Service Gap}
\label{sec:pheromone}
\noindent

\subsubsection*{Stigmergy and environmental memory}
Stigmergy is a form of indirect coordination in which agents communicate
through modifications to a shared environment rather than through direct
interaction \citep{Dorigo1996}. In classical ACO, individual ants deposit chemical pheromones on paths they have traversed; subsequent ants are attracted to stronger trails, reinforcing high-quality routes without any agent needing to communicate directly with another. This framework applies the same principle to transit optimization: the pheromone matrix functions as the algorithm's environmental memory, allowing geographically separated route evaluations conducted in different generations to influence one another through a shared
spatial record.

Concretely, each road segment is mapped to a spatial key derived from its start and end coordinates. Route segments that occupy the same physical corridor therefore contribute to the same pheromone cell even if they belong to different candidate chromosomes. This coordinate-keyed pooling is what makes the memory genuinely stigmergic: it is the environment (the shared matrix) that aggregates experience, not any direct exchange between solutions.

\subsubsection*{Collection and update}
The pheromone matrix is collected \emph{post-simulation}: after a full
agent-based simulation run is completed for a candidate chromosome, the recorded journeys of simulated passengers are used to deposit pheromones along every edge those passengers traversed. Edges that supported efficient journeys receive reinforcement in proportion to the inverse of the total journey cost, following the standard ACO deposition rule \citep{Dorigo1996}:

\[
\Delta\tau_{ij} = \frac{Q}{C}
\]

where $Q$ is the deposition constant and $C$ is the total journey cost. All existing pheromones are then decayed by an evaporation factor $\rho$ before the new deposits are applied:

\[
\tau_{ij}(t+1) = (1 - \rho)\,\tau_{ij}(t) + \Delta\tau_{ij}
\]

Evaporation prevents the matrix from accumulating stale demand signals from routes that are no longer competitive. The evaporation-then-deposition cycle preserves the three mathematical invariants of ACO: path-cost proportional deposition, generational decay, and probabilistic biasing of subsequent operators \citep{Dorigo1996, dorigo2004}.

\subsubsection{The Proportional Demand-Service Disparity Gap}
\label{sec:demand_service_gap}

A primitive approach to identifying unmet transit demand is the linear subtraction of physical vehicle supply from accumulated demand memory. However, in an ant-based framework, this creates a mathematical dimensional mismatch: pheromone intensity ($\tau_{ij}$) is an abstract, continuous variable governed by evaporation and heuristic scaling, whereas transit supply ($s_{ij}$) is a discrete fleet count. A direct subtraction ($\tau_{ij} - s_{ij}$) is inherently unstable because the magnitude of the gap will distort as pheromones accumulate over successive generations.

To resolve this, the framework evaluates structural mismatch using a Proportional Demand-Service Disparity Index. Both the stigmergic demand and the fleet supply are normalized into spatial probability distributions, representing an edge's relative share of the total network activity. 

Let $P_{ij}$ denote an edge's share of the global demand memory, and $S_{ij}$ denote its share of the deployed fleet capacity, calculated as follows:

\[
P_{ij} = \frac{\tau_{ij}}{\sum_{(u,v) \in E} \tau_{uv}}, \qquad S_{ij} = \frac{s_{ij}}{\sum_{(u,v) \in E} s_{uv}}
\]

where:
\begin{itemize}
    \setlength{\itemsep}{0pt}
    \setlength{\parskip}{0pt}
    \setlength{\parsep}{0pt}
    \item $\tau_{ij}$ is the accumulated pheromone intensity on the directed edge $(i, j)$, representing the learned passenger demand.
    \item $s_{ij}$ is the current service supply on edge $(i, j)$, measured as the number of active jeepneys covering that segment scaled by their vehicle weight.
    \item $E$ represents the set of all valid transit edges in the road network graph.
    \item $\sum_{(u,v) \in E} \tau_{uv}$ computes the total sum of pheromones across the entire network.
    \item $\sum_{(u,v) \in E} s_{uv}$ computes the total system-wide service supply currently deployed.
\end{itemize}

The Demand-Service Gap ($\Delta_{ij}$) is then defined as the arithmetic difference between these proportional shares:

\[
\Delta_{ij} = P_{ij} - S_{ij}
\]

This proportional disparity metric fundamentally mirrors spatial equity indices utilized in macroscopic transit connectivity analysis. Welch and Mishra \citeyearpar{welch2013measure} define horizontal equity in transit as ``the equal distribution of an attribute (or recourse) among equal members of a population''. By normalizing the metrics into proportional shares, this framework applies that exact principle to the network topology: the test for equity becomes how well the fleet supply is distributed relative to the underlying spatial demand. 

By binding the gap to the interval $[-1, 1]$, the algorithm becomes immune to scalar pheromone inflation. A positive value ($\Delta_{ij} > 0$) strictly identifies an underserved corridor—meaning the edge's share of passenger demand outpaces its share of the fleet. A negative value ($\Delta_{ij} < 0$) indicates oversupply, where redundant routing causes a corridor to monopolize fleet resources. 

This gap is cached after each evaluation cycle so that the Lamarckian local search operators (such as Spatial Attraction and Oversupply Repulsion) can query the normalized disparity without requiring computationally expensive spatial recalculations.

% ────────────────────────────────────────────────────────────
\subsection{Lamarckian Local Search Operators}
\label{sec:local_search}
\noindent The Lamarckian local search phase applies targeted, intra-chromosome mutations to candidate route systems. Rather than operating randomly, these operators are stigmergically guided by the Proportional Demand-Service Gap calculated from the pheromone matrix. By identifying precisely which corridors are underserved or oversupplied, the operators systematically improve the topological efficiency of the routes. 

\subsubsection{Demand-Driven Or-opt Segment Transplant (Attraction)}
\label{sec:or_opt_attraction}
\noindent The framework implements a Demand-Driven Or-opt Segment Transplant. Or-opt is a classical \ac{VRP} local search move that lifts a contiguous segment of vertices and reinserts it elsewhere at the position of minimum additional cost \citep{laporte2002classical}. Applied to transit design, the operator queries the Demand-Service Gap to identify the most underserved corridors, randomly sampling from the top 20\% of positive-gap edges to maintain stochastic diversity. 

The operator evaluates which $k$-edge window of the target route (where $k \in \{1, 2, 3\}$) can be transplanted adjacent to the underserved corridor with the least net increase in total route length. This reduces the pathfinding burden to exactly three A* calls per evaluated candidate: one to fill the hole left by the lifted segment, and two to bridge the segment into its new position. By sampling candidate positions rather than exhaustively enumerating them, the computational complexity is reduced to $O(N)$, drastically accelerating the simulation while maintaining strict stigmergic alignment.

\subsubsection{Pheromone-Guided 2-opt Segment Reversal (Repulsion)}
\label{sec:2opt_repulsion}
\noindent Transit networks often suffer from structural inefficiencies due to excessive route overlap on central corridors. \cite{kepaptsoglou2009transit} explicitly model this overlap as a primary source of inefficiency that must be penalized in route design. 

The framework implements a Pheromone-Guided 2-opt Segment Reversal. The 2-opt heuristic is the canonical tour-improvement operator; in a transit context, it forces a route to adopt an alternative path by reversing the connections between two cut-points without destroying the closed-loop topology \citep{ciaffi2012feeder}. 

The operator targets the corridor with the most negative Demand-Service Gap (indicating severe oversupply). For a route covering this corridor, the algorithm selects cut-pairs at constrained offsets and generates a new A* path between them. The candidate path is then validated post-hoc to ensure the new segment successfully bypasses the overserved edge. This guarantees that fleet resources are pushed outward into peripheral corridors without requiring destructive modifications to the underlying graph state.

\subsubsection{Demand-Aware Tortuosity Pruning (Sliding Window)}
\label{sec:tortuosity_pruning}
\noindent Directness is a foundational criterion for transit network quality. \cite{ceder1986bus} define circuity (or tortuosity) as the ratio of the actual path length to the straight-line distance, establishing its minimization as a primary design objective. Similarly, \cite{baaj1991transit} utilized explicit route-straightening subroutines to optimize transit layouts.

The operator utilizes a fixed-size sliding window, scanning the route in a single $O(N)$ pass. For each window, it computes the circuity ratio ($\kappa = \frac{L_{path}}{\text{utility}}$). 

To prevent this operator from aggressively unwinding the coverage extensions introduced by the Or-opt Attraction operator, a strict gap-immunity rule is enforced. Any window containing an edge with a positive Demand-Service Gap is exempt from pruning. The most tortuous window lacking this immunity is then replaced with a direct A* shortest path, straightening the route and eliminating low-utility mileage.

\subsubsection{Shared Topological Infrastructure: Loop Assembly and Validation}
\label{sec:loop_validation}
\noindent The framework extracts the surviving, unmodified segments of a parent route and passes them to a dedicated loop builder. This module autonomously bridges any spatial gaps using A* shortest-path routing and explicitly closes the sequence from the tail node back to the head node. Before the candidate route is committed to the simulation state, it must pass a strict closed-loop validation check. If the topological closure fails, the operator safely aborts and returns a null value, ensuring that the original route remains perfectly intact and the multi-layer graph remains free of corruption.

% ────────────────────────────────────────────────────────────\subsection{Memetic Genetic Operators}
\label{sec:genetic}
\noindent
The genetic operators work at the chromosome level, modifying whole jeepney route systems rather than isolated edges. The objective is to improve the entire network as a coordinated service system; crossover and inheritance therefore preserve the strongest structural features of parent networks while opening the search to new route combinations.

Embedding genetic operators directly into an ACO-driven framework represents a deliberate departure from standard ACO mechanics, but it is a scientifically validated methodology  for complex routing problems. As demonstrated by Zhao et al.\ \citeyearpar{zhao2026traversal}, deeply integrating \ac{GA} crossover and mutation operators into the main iterative loop, rather than utilizing them merely as disconnected post-processing tools, acts as a required co-evolutionary mechanism. This hybrid integration dynamically sustains population diversity, counteracts the positive feedback traps inherent to pure pheromone accumulation, and significantly enhances the algorithm's capability to escape local optima without sacrificing the spatial intelligence gathered by the ants.

Two operators are applied in sequence during each reproductive event. Topological Hub Crossover (Section~\ref{sec:crossover}) produces the offspring'sroute geometry by merging the trunk of one parent with complementary feeders from the other. Epigenetic Inheritance (Section~\ref{sec:epigenetic}) then equips the offspring with a pheromone surface that is a fitness-weighted blend of both parents' accumulated demand maps, giving the offspring's subsequent local search an informed prior over which corridors are likely to yield efficient service. The academic rationale for each operator is detailed in the subsections below.

\subsubsection{Topological Hub Crossover (Trunk and Feeder Preservation)}
\label{sec:crossover}

Standard genetic crossover randomly splices route arrays, which can blindly destroy high-performing transit corridors. Gschwender, Jara-Díaz, and Bravo \citeyearpar{gschwender2016feeder} demonstrated mathematically that feeder-trunk network hierarchies outperform flat route systems in urban areas because trunk corridors exhibit economies of density: average operating costs decrease as passenger volumes along a fixed corridor increase. Risso, Nesmachnow, and Faller \citeyearpar{risso2023backbone} extended this by showing that evolutionary algorithms can be explicitly engineered to protect a ``backbone'' trunk network while simultaneously optimizing a broader, peripheral feeder layer, achieving reductions in total system travel time that flat crossover operators cannot match. Topological Hub Crossover implements this principle directly: the crossover preserves the economies-of-density trunk identified through pheromone accumulation, then hybridizes the peripheral coverage layer from the second parent.

The algorithm first identifies the topological hub: the set of road segments with the strongest combined pheromone support and structural importance, representing the busiest corridors where economies of density are highest. Routes from the fitter parent that heavily intersect this hub become the trunk of the child network. Routes from the other parent that do not geometrically conflict with the established trunk are then selected as feeder routes and merged in. The result is a constrained recombination that respects closed-loop validity, route isolation, and corridor utility. The child inherits a coherent trunk-and-feeder structure rather than a fragmented set of edges, mirroring the hierarchical transit design principles established by Gschwender et al.\ \citeyearpar{gschwender2016feeder} and operationalized by Risso et al.\ \citeyearpar{risso2023backbone}.

\subsubsection{Epigenetic Inheritance (Fitness-Weighted Pheromones)}
\label{sec:epigenetic}
\noindent
Offspring produced by crossover inherit route geometry, but without an accompanying demand map, they must rediscover from scratch which corridors support efficient service. Middendorf, Reischle, and Schmeck \citeyearpar{middendorf2002multicolony} proved that in multi-colony ant algorithms, allowing colonies to exchange information about their successful paths prevents stagnation and accelerates convergence toward the global optimum. The mechanism here draws on the same principle, utilizing an approach formalized in evolutionary computation as \textit{Pheromone Crossover} \citep{dong2009multicolony}. Rather than exchanging information between parallel colonies during runtime, the framework transfers accumulated demand memory directly from parents to offspring at the moment of reproduction. Reynolds \citeyearpar{reynolds1994cultural} formalized this as a Cultural Algorithm in which evolution operates simultaneously on a population space (physical routes) and a belief space (accumulated cultural memory). The pheromone matrix is the belief space in this framework; Epigenetic Inheritance is the mechanism by which it propagates.

The child inherits a temporary, localized pheromone matrix that shaped the parents' search history, a methodology structurally validated by Ant-Based Crossover frameworks \citep{barz2003antbased}. If $C_A$ and $C_B$ are the parent costs, the child receives a fitness-weighted pheromone blend derived through continuous arithmetic crossover:

$$ \tau^{\text{child}}_{ij} = w_A \tau^A_{ij} + w_B \tau^B_{ij} $$

with weights:

$$ w_A = \frac{C_B}{C_A + C_B}, \qquad w_B = \frac{C_A}{C_A + C_B} $$

Because the objective is to minimize total system cost, these crossover weights are intentionally inverted. By assigning Parent A the ratio of Parent B's cost over the total, the mathematical formulation guarantees that the parent with the lower (superior) cost receives the proportionately larger weight. This mechanism rigorously enforces the evolutionary balance between exploitation and exploration: the offspring heavily exploits the highly accurate demand map of the superior parent to guide subsequent local searches, while retaining a fractional trace of the secondary parent's map to preserve exploratory diversity and prevent premature convergence.

Critically, this epigenetic transfer does not violate the core principles of Ant Colony Optimization; rather, it functions as an intelligent bootstrap mechanism. In a standard, isolated ACO implementation, the pheromone matrix initializes uniformly with an arbitrary baseline value ($\tau_0$). Because this memetic framework evaluates entirely new offspring topologies every generation, reinitializing to a blank $\tau_0$ would force the simulated passenger agents to blindly rediscover the city's structural demand patterns from scratch, causing severe computational waste.

This redundancy is explicitly addressed by \ac{PACO} architectures. In \ac{PACO}, the pheromone matrix is not a slowly evaporating historical record; rather, all pheromone information corresponds directly to solutions that are members of the actual population \citep{guntsch2002population}. Whenever a solution enters or leaves the population, its influence is explicitly added to or removed from the pheromone matrix \citep{garcia2007taxonomy, guntsch2002population}. By doing so, the ACO construction process acts similarly to Estimation of Distribution Algorithms, allowing the framework to generate new routing solutions based strictly on the distribution of individuals in the current population \citep{garcia2007taxonomy}.

Furthermore, the use of a weighted arithmetic combination to merge distinct pheromone matrices ($\tau^{\text{child}}_{ij} = w_A \tau^A_{ij} + w_B \tau^B_{ij}$) is structurally validated by multi-objective ACO methodologies. For instance, the Pareto Ant Colony Optimization (P-ACO) algorithm utilizes multiple distinct pheromone matrices and combines them at each iteration using a set of computed weights, $p = (p_1, \dots, p_K)$, to evaluate the aggregated pheromone trail via $\sum_{k=1}^K p_k \cdot \tau_{ij}^k$ \citep{garcia2007taxonomy}. By applying this exact weighted combination principle to the matrices of Parent A and Parent B, the framework successfully crosses over multiple heuristic maps without degrading the underlying traffic data.
% ────────────────────────────────────────────────────────────
\section{Optimizer Orchestration and Adaptive Control}
\noindent
The optimizer orchestration layer turns the memetic operators into a reproducible search process. It controls when populations are created, when they are evaluated, when local search is allowed to intervene, and when the search should stop.

% ────────────────────────────────────────────────────────────
\subsection{The Generational Memetic Pipeline}
\label{sec:pipeline}
\noindent
Each generation follows a strict operational sequence so that the optimizer remains deterministic, computationally efficient, and behaviorally interpretable \citep{sastry2013genetic}. The evolutionary lifecycle of a chromosome is orchestrated to minimize the number of times the expensive agent-based temporal simulation must be invoked. 

The pipeline executes the following sequence:
\begin{enumerate}
    \item \textbf{Parent Evaluation:} The pipeline begins by evaluating the current population via the full agent-based temporal simulation. This yields authoritative simulation fitness scores and generates accurate, highly localized pheromone maps for each parent chromosome based on the demand-service gaps experienced by the simulated passenger agents.
    \item \textbf{Selection:} Elitism is applied to copy the best chromosomes forward unchanged. Tournament selection \citep{goldberg1991} then identifies parent pairs for recombination.
    \item \textbf{Crossover and Inheritance:} Topological Hub Crossover produces the child's route geometry. Simultaneously, through Epigenetic Inheritance, the child instantly receives a fitness-weighted blend of the parents' pheromone maps. This warm-start mechanism bypasses the need to run an expensive pre-simulation on the child just to discover where the topological weaknesses lie.
    \item \textbf{Lamarckian Local Search:} The mutation gate decides whether local search should intervene, utilizing the current mutation rate as the trigger probability. Guided by the inherited pheromone map, the mutator attempts structural tweaks. These tweaks are evaluated rapidly using the static Surrogate Cost as a lightweight acceptance baseline. The mutation is kept only if the surrogate cost improves; otherwise, the original geometry is restored.
    \item \textbf{Final Evaluation:} Once the local search is resolved, the fully optimized child undergoes the full agent-based temporal simulation for the first and only time. This single simulation pass calculates the child's official fitness and generates a brand-new pheromone map based on its mutated routes, which it will pass down to future generations.
\end{enumerate}

This sequencing is computationally critical. By evaluating the lightweight surrogate cost during Lamarckian mutation and reserving the temporal simulation strictly for post-mutation evaluation, the framework limits the most expensive operation—the agent-based physics engine—to exactly one pass per child. If local search required full simulations to test every minor structural tweak, or if epigenetic inheritance did not exist to map demand proactively, the evolutionary algorithm would bottleneck immediately. The pipeline therefore preserves optimization correctness, temporal efficiency, and scientific traceability.

% ────────────────────────────────────────────────────────────
\subsection{Adaptive Mutation Scaling via Stagnation Feedback}\label{sec:adaptive}\noindent A static mutation rate creates a fundamental trade-off that cannot be resolved at compile time. Set high, the algorithm explores aggressively but wastes cycles on random perturbations even after a promising region has been found. Set low, the algorithm exploits efficiently but becomes permanently trapped at the first local optimum it encounters. Eiben, Hinterding, and Michalewicz \citeyearpar{eiben1999} formally identified this trade-off and classified stagnation-driven mutation scaling as one of the most effective adaptive mechanisms for escaping local optima in combinatorial optimization. Rather than utilizing a single monolithic mutation rate, the framework deploys three distinct structural mutation operators. To preserve their proportional disruptiveness, each operator is assigned its own independent base probability and maximum cap. The adaptive controller monitors the stagnation counter $s$---the number of consecutive generations without improvement---and scales the activation probability for each specific operator $k$ quadratically toward its respective hard cap:$$ P_{\text{mut}, k} = P_{\text{base}, k} + (P_{\text{max}, k} - P_{\text{base}, k}) \left(\frac{s}{S_{\text{limit}}}\right)^2 $$This explicit decoupling ensures that while all three operators follow the same quadratic escalation schedule during stagnation, their individual activation thresholds remain structurally balanced. The quadratic curve delibrerately provides patience early in stagnation, when natural escape is still probable, and applies force sharply as the plateau persists \citep{abulail2025, farda2024}. A linear schedule would apply uniform pressure regardless of how severe the stagnation is, penalizing early generations unnecessarily and responding too slowly near the limit. The moment a new best solution is found, the stagnation counter resets to zero and all three mutation rates return to their respective baselines, preventing the controller from continuing to force exploration after the algorithm has already broken free. The same controller also governs the Lamarckian local search probability and its spatial intensity. Unlike the stagnation-driven mutation rates, these local search parameters decay linearly across generations:$$ P_{\text{local}}(g) = P_{\min} + (P_{\max} - P_{\min}) \left(1 - \frac{g}{G_{\max}}\right) $$Early generations apply broad route edits to rapidly shape candidate topologies; later generations make conservative, surgical refinements. This coupling prevents the search from being too disruptive near convergence \citep{eiben1999}.

% ────────────────────────────────────────────────────────────
\subsection{Multi-Dimensional Convergence Criteria (Jaccard and Variance)}
\label{sec:convergence}
\noindent
The optimizer stops only when both structural and numerical evidence indicate that the search has exhausted useful diversity. Two independent signals are monitored; Section~\ref{sec:metrics} explains the mathematical properties of both metrics in detail.

The first criterion is elite Jaccard similarity, which measures convergence in the \textit{decision space} (phenotypic saturation). If $E_A$ and $E_B$ are the active edge sets of two elite chromosomes, their structural overlap is:

$$ J(A,B) = \frac{|E_A \cap E_B|}{|E_A \cup E_B|} $$

The optimizer computes the average pairwise Jaccard similarity among the top ten percent of the population. When that value remains above the configured threshold for a sufficient number of consecutive generations (\texttt{jaccard\_patience}), the population is considered phenotypically saturated: the best route systems have become structurally too similar to recombine meaningfully \citep{sastry2013genetic}.

The second criterion is fitness variance, which measures convergence in the \textit{objective space}. The algorithm tracks the dispersion of fitness scores across the current population ($P$) of size $N$, where $\mu$ is the mean population fitness:

$$ \sigma^2(F) = \frac{1}{N} \sum_{i=1}^{N} (F_i - \mu)^2 $$

When $\sigma^2(F)$ collapses below a minimal tolerance threshold ($\epsilon$) for a sustained number of generations, the objective gradient is considered exhausted. This criterion is theoretically vital because transit routing frequently exhibits "neutral networks" or flat fitness landscapes. In such scenarios, multiple chromosomes might differ by a few parallel streets—keeping their Jaccard similarity low enough to evade the first stopping criterion—but these minor topological variations yield the exact same generalized travel time for passengers. 

If the optimizer relied solely on Jaccard similarity, it would waste massive amounts of computation endlessly mutating routes on a flat plateau where no numerical improvement is possible. By enforcing both topological (Jaccard) and objective (variance) dispersion checks, the optimizer strictly avoids premature termination while guaranteeing that it does not loop infinitely when the numerical search gradient has vanished.

% ────────────────────────────────────────────────────────────
\subsection{Telemetry, Deterministic Replay, and State Preservation}
\label{sec:telemetry}
\noindent
Telemetry records how the optimizer behaves over time, while state preservation ensures that the same behavior can be resumed exactly after an interruption. The telemetry engine writes three complementary outputs. The history log stores generational fitness statistics: best cost, mean cost, mutation rate, and stagnation count. The lineage log records parent-child relationships so that any final chromosome can be traced back through its full ancestry. The \ac{JSON} snapshot stream exports route geometries, pheromone hotspots, and demand-choked corridors for downstream visualization.

State preservation is implemented through atomic checkpointing. The live optimization state—population, pheromone matrices, generation counter, stagnation counter, best fitness, and pseudorandom number generator state—is first serialized to a temporary file, then promoted to a permanent checkpoint only after the write completes successfully. This prevents corrupted partial saves. Deterministic replay is maintained by restoring the random state only after all deterministic infrastructure has been rebuilt, so the resumed run consumes exactly the same sequence of random numbers that the uninterrupted run would have used.

% ────────────────────────────────────────────────────────────
\section{Post-Optimization Evaluation Metrics}
\label{sec:metrics}
\noindent
After optimization, the final route systems are evaluated using several complementary metrics. No single score is sufficient because a network can be similar in structure, similar in geometry, similar in coverage distribution, or similar in ranking behavior without being identical in all four senses. The metrics documented here are implemented as pure, stateless functions in \texttt{evaluation\_metrics.py} and applied to domain objects through \texttt{post\_evaluation.py}. This two-layer separation ensures that each metric is independently testable and that its mathematical assumptions are transparent.

% ────────────────────────────────────────────────────────────
\subsection{Topological and Geometric Similarity (Jaccard, Discrete Fréchet, GED)}
\label{sec:topo_metrics}
\noindent

\paragraph{Jaccard similarity.}
Topological similarity is measured with the Jaccard index over active route edges. For two edge sets $A$ and $B$:

\[
J(A,B) = \frac{|A \cap B|}{|A \cup B|}
\]

If system $A$ and system $B$ each use 100 street segments and 80 of those are identical, $J = 80/120 = 0.667$. The metric is symmetric, bounded in $[0,1]$, and requires no parameter tuning. It answers the question of how much corridor overlap exists between two route systems—whether two candidate networks share the same trunks and feeder branches. The Jaccard index appears in three distinct roles throughout this framework: inside the optimizer for convergence detection (Section~\ref{sec:convergence}), in Topological Hub Crossover for trunk identification (Section~\ref{sec:crossover}), and here in post-evaluation for system comparison. A single implementation in \texttt{evaluation\_metrics.py} serves all three use cases.

\paragraph{Discrete Fréchet distance.}
Jaccard similarity does not capture the ordering of nodes along a route, so it is complemented by the discrete Fréchet distance \citep{eiter1994computing}. The Fréchet distance is often described as the ``dog-walker'' metric: a person walks along path $P$ and a dog along path $Q$, connected by a leash; neither may backtrack. The Fréchet distance is the minimum leash length that allows both to complete their traversals. The discrete variant restricts the comparison to the vertex sequences of each path and is computed via a standard $O(n \times m)$ dynamic programming recurrence:

\[
ca(i,j) = \max\!\bigl(\min(ca(i{-}1,j),\, ca(i,j{-}1),\, ca(i{-}1,j{-}1)),\,
             d(P_i, Q_j)\bigr)
\]

where $d(P_i, Q_j)$ is the distance between the $i$-th vertex of $P$ and the $j$-th vertex of $Q$. Critically, Fréchet distance respects traversal order: two routes sharing identical streets but traveling in opposite directions correctly register as dissimilar, making it strictly more informative than Hausdorff distance for evaluating route mutations.

\paragraph{Graph Edit Distance (GED).}
Graph Edit Distance (GED) provides a third, strictly topological perspective by calculating the minimum cost required to morph the structure of one route network into another \citep{sanfeliu1983distance}. Formally, given two graphs $G_1$ and $G_2$, the distance is defined as the minimum cost of an edit path $P$ consisting of a sequence of operations $o \in \{ \text{insertion, deletion, substitution} \}$ such that:

$$ d(G_1, G_2) = \min_{P \in \mathcal{P}} \sum_{o \in P} c(o) $$

where $c(o)$ represents the cost associated with each transformation \citep{sanfeliu1983distance}. While Jaccard similarity evaluates set-theoretic edge overlap and Fréchet distance measures continuous spatial trajectories, both fail to capture the discrete sequence of network connectivity. For example, two transit routes might run spatially parallel (yielding a low Fréchet distance) or share a large pool of unordered road segments (yielding a high Jaccard score) while possessing entirely different intersection sequences and transfer structures. GED resolves this blind spot by quantifying explicit topological operations: node insertions (adding new stops), node deletions (removing service from a point), and edge substitutions (rerouting the connective topology). Taken together, Jaccard, Fréchet, and GED separate the evaluation of route diversity into three orthogonal mathematical dimensions: shared corridor membership (set theory), continuous route shape (geometry), and structural connectivity (topology). This multi-dimensional evaluation ensures that the optimization framework can accurately distinguish between visually similar but operationally distinct transport networks.

% ────────────────────────────────────────────────────────────
\subsection{Distributional Distance and Spatial Coverage (1D/2D Wasserstein)}
\label{sec:wasserstein}
\noindent
Wasserstein distance, also referred to as the Earth Mover's Distance, measures how much mass must be moved to transform one distribution into the shape of another, at a cost equal to mass times distance moved \citep{debacco2023optimaltransport}. Formally, for two probability distributions $\mu$ and $\nu$ over a metric space:

\[
W_1(\mu,\nu) = \inf_{\gamma \in \Gamma(\mu,\nu)}
  \int_{X \times X} d(x,y)\, d\gamma(x,y)
\]

where $\Gamma(\mu,\nu)$ is the set of all joint distributions with marginals
$\mu$ and $\nu$.

Two variants are implemented. The one-dimensional form
(\texttt{scipy.stats.wasserstein\_distance}) compares univariate empirical distributions—commute times, route lengths, or waiting times—and evaluates whether the final route system shifts travel outcomes toward lower costs. The two-dimensional form applies optimal transport to weighted spatial point distributions, comparing where service is concentrated on the map and evaluating whether coverage has migrated toward underserved spatial regions.

Unlike Jaccard, which is purely topological, and Fréchet, which is purely geometric, Wasserstein is demand-aware. Two systems covering different streets but serving the same demand surface produce a low Wasserstein distance; two systems covering the same streets but serving different demand concentrations produce a high one. De Bacco et al.\ \citeyearpar{debacco2023optimaltransport} applied optimal transport theory to compare multi-layer urban network structures in a formulation directly analogous to comparing the demand coverage distributions of two candidate jeepney route systems.

% ────────────────────────────────────────────────────────────
\subsection{Network Diversity (Shannon Entropy)}
\label{sec:entropy}
\noindent
Shannon entropy measures how evenly route usage is distributed across the network. If $p_i$ is the normalized share of usage assigned to route $i$:

\[
H = -\sum_i p_i \log p_i
\]

High entropy indicates that traffic and service are spread across many routes, suggesting genuine multi-modal choice. Low entropy indicates concentration on a small number of routes, which may be acceptable in an intentionally hierarchical network but may also signal premature convergence or over-reliance on a narrow corridor set. Levinson \citeyearpar{levinson2012networkstructure} applied entropy measures specifically to urban transportation networks and demonstrated that structural diversity correlates with network resilience: systems with higher routing entropy are more robust to link failures and demand perturbations. Entropy therefore complements the similarity metrics by describing network breadth rather than overlap.

% ────────────────────────────────────────────────────────────
\subsection{Surrogate Fidelity Validation (Spearman and Kendall Rank Correlation)}
\label{sec:surrogate_val}
\noindent
The surrogate evaluator is useful only insofar as it preserves the ordering of route systems produced by the full simulation. Evolutionary selection is ranking-based: tournament selection picks the fittest of $k$ chromosomes, and elitism preserves the top $n$. These operations depend on relative ordering, not absolute magnitude. Jin \citeyearpar{jin2005survey} explicitly established Spearman rank correlation as the primary validation metric for fitness approximation in evolutionary computation, noting that a surrogate's value lies in its ability to rank solutions correctly, not in reproducing their exact numerical scores.

To test this, surrogate scores and full simulation scores are compared using two rank-based coefficients. If $d_i$ is the rank difference for chromosome $i$, Spearman correlation is:

\[
\rho = 1 - \frac{6 \sum d_i^2}{n(n^2 - 1)}
\]

where $n$ is the number of chromosomes in the validation set. Kendall
correlation—more robust for small sample sizes—is:

\[
\tau = \frac{C - D}{\binom{n}{2}}
\]

where $C$ is the number of concordant pairs and $D$ is the number of discordant pairs. A Spearman $\rho \geq 0.7$ indicates sufficient ranking fidelity to guide local search reliably; below $0.5$, the surrogate would actively mislead the selection pressure \citep{jin2005survey}. The framework additionally reports Kendall $\tau$, normalized \ac{RMSE}, and \ac{MAPE} as supporting diagnostics, but Spearman $\rho$ is the decision-making metric.